

\section{The OSI 7 layer model }

\begin{description}



\os{newpage} item[{\rot\bf  Structure }]~\\[-3ex]

\bi
  \ite General

  \ite The 7 layers

  \ite Their tasks listed individually

  \ite Communication between the layers

  \ite Sources

\ei


\os{newpage} item[{\rot\bf  General }]~\\[-3ex]

\bi
  \ite Developed since 1979 \& standardised by ISO (International Organisation for Standardisation)

  \ite It consists of individual open layers

  \ite Serves the communication of information-processing systems as the basis for a number of manufacturerindependent network protocols.

\ei


\os{newpage} item[{\rot\bf  General }]~\\[-3ex]

\bi
  \ite Describes a unified procedure and rules for the exchange of data

  \ite Finds application in public communication technology

  \ite It is more or less a translation programme for data transfer between different systems.

\ei


\os{newpage} item[{\rot\bf  General }]~\\[-3ex]

\bi
  \ite Is hierarchically structured

  \ite Lower layer provides its services to the one above it

\ei


\os{newpage} item[{\rot\bf  The 7 layers }]~\\[-3ex]

\bi
  \ite From top to bottom

  \ite
  \begin{enumerate}
    \setcounter{enumii}{6}
    \ite application layer
  \end{enumerate}
  \ite
  \begin{enumerate}
    \setcounter{enumii}{5}
    \ite presentation layer
  \end{enumerate}
  \ite
  \begin{enumerate}
    \setcounter{enumii}{4}
    \ite session layer
  \end{enumerate}
  \ite
  \begin{enumerate}
    \setcounter{enumii}{3}
    \ite transport layer
  \end{enumerate}
  \ite
  \begin{enumerate}
    \setcounter{enumii}{2}
    \ite network layer (network layer)
  \end{enumerate}
  \ite
  \begin{enumerate}
    \setcounter{enumii}{1}
    \ite data link layer
  \end{enumerate}
  \ite
  \begin{enumerate}
    \ite physical layer (bit transmission layer)
  \end{enumerate}
\ei


\os{newpage} item[{\rot\bf  The 7 layers }]~\\[-3ex]

\bi
  \ite Layers 1-4 are assigned to the transport function

  \ite Layers 5 - 7 are assigned to the user functions

\ei



\os{newpage} item[{\rot\bf  The OSI Model (Open Systems Interconnection) }]~\\[-3ex]

\begin{center}
\includegraphics[max width=\textwidth]{2023_01_13_b07e3b4bcffdaab9792bg-009}
\end{center}

\begin{center}
\includegraphics[max width=\textwidth]{2023_01_13_b07e3b4bcffdaab9792bg-009(1)}
\end{center}

\begin{center}
\includegraphics[max width=\textwidth]{2023_01_13_b07e3b4bcffdaab9792bg-009(2)}
\end{center}

(c) Copyright 2008 Steven Iveson WWW. \href{http://networkstuff.eu}{networkstuff.eu}

\begin{center}
\includegraphics[max width=\textwidth]{2023_01_13_b07e3b4bcffdaab9792bg-009(3)}
\end{center}


\os{newpage} item[{\rot\bf  Reference Models }]~\\[-3ex]

Layer

Name of unit exchanged

The OSI reference model.

\begin{center}
\includegraphics[max width=\textwidth]{2023_01_13_b07e3b4bcffdaab9792bg-010}
\end{center}


\os{newpage} item[{\rot\bf  Layer 7 }]~\\[-3ex]

\bi
  \ite Is the top layer in the system

  \ite establishes the connection between application programmes (e.g.: email, data transfer, etc.)

  \ite Basic services:

\ei
\bi
 \ite File exchange
 \ei


\bi
  \ite Management protocols for user access, file access rights, electronic mail, etc.
\ei


\os{newpage} item[{\rot\bf  Layer 6 }]~\\[-3ex]

\bi
  \ite It provides functions for protocol conversion, data transmission and data encryption

  \ite Takes over encoding \& decoding (e.g. EBCDIC in ASCII)

  \ite Takes over the definition of the formats and control characters for the exchange of the data

\ei


\os{newpage} item[{\rot\bf  Layer 6 }]~\\[-3ex]


\os{newpage} item[{\rot\bf  - Takes care of the syntax \& correctness of the data }]~\\[-3ex]

\os{newpage} item[{\rot\bf Layer 5 }]~\\[-3ex]

\os{newpage} item[{\rot\bf  - Controls the establishment, execution and termination of the connection }]~\\[-3ex]

\textbackslash author\{

\bi
  \ite Takes over monitoring of operating parameters, data \textbackslash  flow control, reconnection in case of error and \textbackslash  synchronisation between user actions
\}
\ei


\os{newpage} item[{\rot\bf  $\underline{\text { Layer } 4}$ }]~\\[-3ex]


\os{newpage} item[{\rot\bf  - Is the connection between application-oriented \& transport-oriented layers }]~\\[-3ex]

\bi
  \ite Ensures that all data packets reach the correct
recipient
\ei


\os{newpage} item[{\rot\bf  $\underline{\operatorname{Layer}} 4$ }]~\\[-3ex]

\bi
  \ite Establishing the data connection between two partners, data transport, flow control, error detection and correction

  \ite Ensuring completeness, uniqueness \& sequence of data blocks

  \ite Their services are divided into 4 classes

\ei


\os{newpage} item[{\rot\bf  $\underline{\operatorname{Layer} 4}$ }]~\\[-3ex]


\os{newpage} item[{\rot\bf  - Class O: }]~\\[-3ex]


\os{newpage} item[{\rot\bf  o Is the simplest }]~\\[-3ex]
\bi
 \ite No error handling
 \ei

\bi
 \ite there is no error control compared to layer 3
 \ei

\bi
 \ite One transport connection corresponds to exactly one network connection
 \ei



\os{newpage} item[{\rot\bf  Layer 4 }]~\\[-3ex]


\os{newpage} item[{\rot\bf  - Class 1: }]~\\[-3ex]
\bi
 \ite No error handling
 \ei

\bi
 \ite Attempts to correct errors reported by layer 3 and does not forward them to layer 5 .
 \ei

\bi
 \ite For example, if the transport connection is interrupted, an attempt can be made to rebuild it without this being noticed above layer 4 .
 \ei



\os{newpage} item[{\rot\bf  Layer 4 }]~\\[-3ex]


\os{newpage} item[{\rot\bf  - Class 2: }]~\\[-3ex]


\os{newpage} item[{\rot\bf  o can establish several transport connections (multiplex connection) }]~\\[-3ex]


\os{newpage} item[{\rot\bf  o In this case, the network connection may only be disconnected when the last transport connection has been terminated. }]~\\[-3ex]


\os{newpage} item[{\rot\bf  $\underline{\operatorname{Layer} 4}$ }]~\\[-3ex]


\os{newpage} item[{\rot\bf  - Class 3 : }]~\\[-3ex]


\os{newpage} item[{\rot\bf  o Covers the services of classes 1 \& 2 (error handling \& multiplex connections) }]~\\[-3ex]


\os{newpage} item[{\rot\bf  Layer 4 }]~\\[-3ex]

\bi
  \ite Class 4:
\ei
\bi
 \ite Contains, in addition to the functions of class 3, additional mechanisms for error detection/handling
 \ei


\bi
  \ite Especially in datagram-oriented networks (LAN), a connection-oriented service can be provided in this way (ensuring completeness, uniqueness and sequence of data blocks).
\ei


\os{newpage} item[{\rot\bf  Layer 4 }]~\\[-3ex]

\bi
  \ite Devices:
\ei
\bi
 \ite Gateway
 \ei



\os{newpage} item[{\rot\bf  Layer 3 }]~\\[-3ex]

\bi
  \ite is mainly used for data packet
transmission

  \ite Responsible for the choice of data
paths (routing)

  \ite for multiplexing several connections over individual sections

\ei


\os{newpage} item[{\rot\bf  Layer 3 }]~\\[-3ex]

\bi
  \ite For error handling, flow control \& flow control between the endpoints of a connection (not between the user processes).

  \ite Removal of bottlenecks (consequence of data
collisions)

\ei


\os{newpage} item[{\rot\bf  Layer 3 }]~\\[-3ex]

\bi
  \ite Devices:
\ei
\bi
 \ite Router
 \ei

\bi
 \ite Layer 3 Switch
 \ei



\os{newpage} item[{\rot\bf  Layer 2 }]~\\[-3ex]

\bi
  \ite Ensuring a functioning connection between two directly adjacent stations

  \ite Provides a defined framework for data transport, error detection and synchronisation of data

\ei


\os{newpage} item[{\rot\bf  Layer 2 }]~\\[-3ex]

\bi
  \ite Information is divided into blocks of appropriate length called frames and provided with check info for error detection/correction.

  \ite Flow control for the binary data

  \ite Transmitter \& receiver hardware addresses are added to the data frame

\ei


\os{newpage} item[{\rot\bf  Layer 2 }]~\\[-3ex]

\bi
  \ite For local networks, this layer is subdivided again: - 2 a (Media Access Control, MAC): Controls access to the transmission medium
\ei
\bi
 \ite 2 b (Logical Link Control, LLC): Layer 2 functions dependent on the transmission medium
 \ei



\os{newpage} item[{\rot\bf  Layer 2 }]~\\[-3ex]

\bi
  \ite Devices:
\ei
\bi
 \ite Bridge
 \ei


Switch


\os{newpage} item[{\rot\bf  Layer 1 }]~\\[-3ex]

\bi
  \ite Is the bottom layer

  \ite This is where the physical data transmission takes place

  \ite It defines the electrical, mechanical \& functional parameters for the physical connection of two units (e.g. level, modulation, cable, connector, transmission rate, etc.).

\ei


\os{newpage} item[{\rot\bf  Layer 1 }]~\\[-3ex]

\bi
  \ite Devices and network components:
\ei
\bi
 \ite Antenna
 \ei

\bi
 \ite Amplifier
 \ei

\bi
 \ite Plug and socket for the mains cable
 \ei

\bi
 \ite Repeater
 \ei

\bi
 \ite $\mathrm{Hub}$
 \ei

\bi
 \ite Transceiver
 \ei

\bi
 \ite T-piece and the end resistor
 \ei



\os{newpage} item[{\rot\bf  Communication between the layers }]~\\[-3ex]

Layer 7

Layer 6

\begin{center}
\includegraphics[max width=\textwidth]{2023_01_13_b07e3b4bcffdaab9792bg-032}
\end{center}

Layer 5

Layer 4

\begin{center}
\begin{tabular}{|c|c|}
\hline
TCP Header & TCP Data \\
\hline
MAC Header & IP Data \\
\hline
Ethernet/ Token-Ring Frame / X.25 / FR / ATM Data &  \\
\hline
\end{tabular}
\end{center}

Layer 3

Layer 1

Layer 2 TCP Frame

Application

Frame

IP Frame

MAC Frame


\os{newpage} item[{\rot\bf  Communication between the layers }]~\\[-3ex]

\begin{center}
\includegraphics[max width=\textwidth]{2023_01_13_b07e3b4bcffdaab9792bg-033}
\end{center}

Figure 2. Encapsulation and Decapsulation of TCP/IP Protocol


\os{newpage} item[{\rot\bf  Devices }]~\\[-3ex]

OSI (Open Source Interconnection) 7 Layer Model

\begin{center}
\includegraphics[max width=\textwidth]{2023_01_13_b07e3b4bcffdaab9792bg-034}
\end{center}


\os{newpage} item[{\rot\bf  Sources }]~\\[-3ex]

\bi
  \ite \href{http://www.wikipedia.de}{www.wikipedia.de} 17.12.05

  \ite \href{http://homepages.luc.edu/}{http://homepages.luc.edu/} bmontes/CIEP489Images/OSI-Model.gif 17.12.05

  \ite \href{http://www.netzmafia.de/skripten/netze/netzo.htm}{http://www.netzmafia.de/skripten/netze/netzo.htm} l\#0.1 17.12.05

  \ite My DVT binder 2004/05

\ei


\os{newpage} item[{\rot\bf  Transceiver Block Diagram }]~\\[-3ex]


\os{newpage} item[{\rot\bf  Transmitter }]~\\[-3ex]

Signal

\begin{center}
\includegraphics[max width=\textwidth]{2023_01_13_b07e3b4bcffdaab9792bg-037}
\end{center}

\begin{center}
\includegraphics[max width=\textwidth]{2023_01_13_b07e3b4bcffdaab9792bg-037(1)}
\end{center}


\os{newpage} item[{\rot\bf  Receiver }]~\\[-3ex]

Antenna


\os{newpage} item[{\rot\bf  Low Noise Amplifier }]~\\[-3ex]

\begin{center}
\includegraphics[max width=\textwidth]{2023_01_13_b07e3b4bcffdaab9792bg-038}
\end{center}

$\square$ Low Noise Amplifier

Gain-Stage Amplifier


\os{newpage} item[{\rot\bf  Transceiver $=$ Transmitter $+$ Receiver }]~\\[-3ex]

\begin{center}
\includegraphics[max width=\textwidth]{2023_01_13_b07e3b4bcffdaab9792bg-039}
\end{center}


\os{newpage} item[{\rot\bf  Signal Encoding Techniques }]~\\[-3ex]


\os{newpage} item[{\rot\bf  Encoding Techniques }]~\\[-3ex]

\bi
  \ite Digital data, digital signal

  \ite Analog data, digital signal

  \ite Digital data, analog signal

  \ite Analog data, analog signal

\ei


\os{newpage} item[{\rot\bf  Digital Data, Digital Signal }]~\\[-3ex]

\bi
  \ite Digital signal

  \ite Discrete, discontinuous voltage pulses

\ei
\bi
 \ite Each pulse is a signal element
 \ei


\bi
  \ite Binary data encoded into signal elements
\ei


\os{newpage} item[{\rot\bf  Terms (1) }]~\\[-3ex]

\bi
  \ite Unipolar
\ei

All signal elements have same sign

\bi
  \ite Polar

  \ite One logic state represented by positive voltage the other by negative voltage

  \ite Data rate

\ei
\bi
 \ite Rate of data transmission in bits per second
 \ei


\bi
  \ite Duration or length of a bit
\ei
\bi
 \ite Time taken for transmitter to emit the bit
 \ei



\os{newpage} item[{\rot\bf  Terms (2) }]~\\[-3ex]

\bi
  \ite Modulation rate

  \ite Rate at which the signal level changes

  \ite Measured in baud = signal elements per second

  \ite Mark and Space

\ei
\bi
 \ite Binary 1 and Binary o respectively
 \ei



\os{newpage} item[{\rot\bf  Interpreting Signals }]~\\[-3ex]

\bi
  \ite Need to know
\ei
\bi
 \ite Timing of bits - when they start and end o Signal levels
 \ei


\bi
  \ite Factors affecting successful interpreting of signals o Signal to noise ratio
\ei
\bi
 \ite Data rate
 \ei

\bi
 \ite Bandwidth
 \ei



\os{newpage} item[{\rot\bf  Comparison of Encoding Schemes (1) }]~\\[-3ex]

\bi
  \ite Signal Spectrum

  \ite Lack of high frequencies reduces required bandwidth

  \ite Lack of dc component allows ac coupling via transformer, providing isolation

\ei
\bi
 \ite Concentrate power in the middle of the bandwidth
 \ei



\os{newpage} item[{\rot\bf  - Clocking }]~\\[-3ex]
\bi
 \ite Synchronizing transmitter and receiver
 \ei

\bi
 \ite External clock
 \ei

\bi
 \ite Sync mechanism based on signal
 \ei



\os{newpage} item[{\rot\bf  Comparison of Encoding Schemes (2) }]~\\[-3ex]

\bi
  \ite Error detection

  \ite Can be built in to signal encoding

  \ite Signal interference and noise immunity

\ei
\bi
 \ite Some codes are better than others
 \ei


\bi
  \ite Cost and complexity

  \ite Higher signal rate (\& thus data rate) lead to higher costs

\ei
\bi
 \ite Some codes require signal rate greater than data rate
 \ei



\os{newpage} item[{\rot\bf  Encoding Schemes }]~\\[-3ex]

\bi
  \ite Nonreturn to Zero-Level (NRZ-L)

  \ite Nonreturn to Zero Inverted (NRZI)

  \ite Bipolar -AMI

  \ite Pseudoternary

  \ite Manchester

  \ite Differential Manchester

  \ite B8ZS

  \ite $\mathrm{HDB}_{3}$

\ei


\os{newpage} item[{\rot\bf  Nonreturn to Zero-Level (NRZ-L) }]~\\[-3ex]

\bi
  \ite Two different voltages for 0 and 1 bits

  \ite Voltage constant during bit interval o no transition I.e. no return to zero voltage

  \ite e.g. Absence of voltage for zero, constant positive voltage for one

  \ite More often, negative voltage for one value and positive for the other

  \ite This is NRZ-L

\ei


\os{newpage} item[{\rot\bf  Nonreturn to Zero Inverted }]~\\[-3ex]

\bi
  \ite Nonreturn to zero inverted on ones

  \ite Constant voltage pulse for duration of bit

  \ite Data encoded as presence or absence of signal transition at beginning of bit time

  \ite Transition (low to high or high to low) denotes a binary 1

  \ite No transition denotes binary o

  \ite An example of differential encoding

\ei

\begin{center}
\includegraphics[max width=\textwidth]{2023_01_13_b07e3b4bcffdaab9792bg-051}
\end{center}


\os{newpage} item[{\rot\bf  Differential Encoding }]~\\[-3ex]

\bi
  \ite Data represented by changes rather than levels

  \ite More reliable detection of transition rather than level

  \ite In complex transmission layouts it is easy to lose sense of polarity

\ei


\os{newpage} item[{\rot\bf  NRZ pros and cons }]~\\[-3ex]

\bi
  \ite Pros
\ei
\bi
 \ite Easy to engineer
 \ei


\bi
  \ite Make good use of bandwidth

  \ite Cons

\ei
\bi
 \ite dc component
 \ei

\bi
 \ite Lack of synchronization capability
 \ei


\bi
  \ite Used for magnetic recording

  \ite Not often used for signal transmission

\ei


\os{newpage} item[{\rot\bf  Multilevel Binary }]~\\[-3ex]

\bi
  \ite Use more than two levels

  \ite Bipolar-AMI

\ei
\bi
 \ite zero represented by no line signal
 \ei

\bi
 \ite one represented by positive or negative pulse
 \ei

\bi
 \ite one pulses alternate in polarity
 \ei


\bi
  \ite No loss of sync if a long string of ones (zeros still a problem)
\ei
\bi
 \ite No net dc component
 \ei

\bi
 \ite Lower bandwidth
 \ei

\bi
 \ite Easy error detection
 \ei



\os{newpage} item[{\rot\bf  Pseudoternary }]~\\[-3ex]

\bi
  \ite One represented by absence of line signal

  \ite Zero represented by alternating positive and negative

  \ite No advantage or disadvantage over bipolar-AMI

\ei


\os{newpage} item[{\rot\bf  Bipolar-AMI and Pseudoternary }]~\\[-3ex]


\os{newpage} item[{\rot\bf  Bipolar-AMI }]~\\[-3ex]

(most recent preceding 1 bit has negative voltage)


\os{newpage} item[{\rot\bf  Pseudoternary }]~\\[-3ex]

(most recent preceding 0 bit has negative voltage)

\begin{center}
\includegraphics[max width=\textwidth]{2023_01_13_b07e3b4bcffdaab9792bg-056}
\end{center}


\os{newpage} item[{\rot\bf  Trade Off for Multilevel Binary }]~\\[-3ex]

\bi
  \ite Not as efficient as NRZ
\ei
\bi
 \ite Each signal element only represents one bit
 \ei

\bi
 \ite In a 3 level system could represent $\log _{2} 3=1.58$ bits
 \ei


\bi
  \ite Receiver must distinguish between three levels $(+\mathrm{A},-\mathrm{A}, \mathrm{o})$

  \ite Requires approx. $3 \mathrm{~dB}$ more signal power for same probability of bit error

\ei


\os{newpage} item[{\rot\bf  Biphase }]~\\[-3ex]

\bi
  \ite Manchester
\ei
\bi
 \ite Transition in middle of each bit period
 \ei

\bi
 \ite Transition serves as clock and data
 \ei

\bi
 \ite Low to high represents one
 \ei

\bi
 \ite High to low represents zero
 \ei

\bi
 \ite Used by IEEE 802.3
 \ei


\bi
  \ite Differential Manchester
\ei
\bi
 \ite Midbit transition is clocking only
 \ei

\bi
 \ite Transition at start of a bit period represents zero
 \ei

\bi
 \ite No transition at start of a bit period represents one
 \ei

\bi
 \ite Note: this is a differential encoding scheme
 \ei

\bi
 \ite Used by IEEE 802.5
 \ei



\os{newpage} item[{\rot\bf  Manchester Encoding }]~\\[-3ex]


\os{newpage} item[{\rot\bf  Manchester Encoding }]~\\[-3ex]

\begin{center}
\includegraphics[max width=\textwidth]{2023_01_13_b07e3b4bcffdaab9792bg-059}
\end{center}


\os{newpage} item[{\rot\bf  Differential Manchester Encoding }]~\\[-3ex]


\os{newpage} item[{\rot\bf  Differential Manchester Encoding }]~\\[-3ex]

\begin{center}
\includegraphics[max width=\textwidth]{2023_01_13_b07e3b4bcffdaab9792bg-060}
\end{center}

\os{newpage} item[{\rot\bf Biphase Pros and Cons }]~\\[-3ex]

\os{newpage} item[{\rot\bf  - Con }]~\\[-3ex]

At least one transition per bit time and possibly two
 - Maximum modulation rate is twice NRZ

 o Requires more bandwidth

\os{newpage} item[{\rot\bf  - Pros }]~\\[-3ex]




\os{newpage} item[{\rot\bf  Modulation Rate }]~\\[-3ex]

5 bits $=5 \mu \mathrm{sec}$

\begin{center}
\includegraphics[max width=\textwidth]{2023_01_13_b07e3b4bcffdaab9792bg-062}
\end{center}

1

1

1

1

1

NRZI
\includegraphics[max width=\textwidth, center]{2023_01_13_b07e3b4bcffdaab9792bg-062(1)}

Manchester

1 bit $=$

1 signal element $=$

1 μ $\mathrm{sec}$

\begin{center}
\includegraphics[max width=\textwidth]{2023_01_13_b07e3b4bcffdaab9792bg-062(2)}
\end{center}

$1 \mathrm{bit}=$

l $\mu$ sec

\begin{center}
\includegraphics[max width=\textwidth]{2023_01_13_b07e3b4bcffdaab9792bg-062(3)}
\end{center}

\begin{center}
\includegraphics[max width=\textwidth]{2023_01_13_b07e3b4bcffdaab9792bg-062(4)}
\end{center}

1

1 signal element $=$

$0.5 \mu \mathrm{sec}$


\os{newpage} item[{\rot\bf  Scrambling }]~\\[-3ex]

\bi
  \ite Use scrambling to replace sequences that would produce constant voltage

  \ite Filling sequence

  \ite Must produce enough transitions to sync

  \ite Must be recognized by receiver and replace with original

  \ite Same length as original

  \ite No dc component

  \ite No long sequences of zero level line signal

  \ite No reduction in data rate

  \ite Error detection capability

\ei


\os{newpage} item[{\rot\bf  B8ZS }]~\\[-3ex]

\bi
  \ite Bipolar With 8 Zeros Substitution

  \ite Based on bipolar-AMI

  \ite If octet of all zeros and last voltage pulse preceding was positive encode as $000+-0-+$

  \ite If octet of all zeros and last voltage pulse preceding was negative encode as 000-+0+-

  \ite Causes two violations of AMI code

  \ite Unlikely to occur as a result of noise

  \ite Receiver detects and interprets as octet of all zeros

\ei


\os{newpage} item[{\rot\bf  $\mathrm{HDB}_{3}$ }]~\\[-3ex]

\bi
  \ite High Density Bipolar 3 Zeros

  \ite Based on bipolar-AMI

  \ite String of four zeros replaced with one or two pulses

\ei


\os{newpage} item[{\rot\bf  B8ZS and $\mathrm{HDB}_{3}$ }]~\\[-3ex]

\begin{center}
\includegraphics[max width=\textwidth]{2023_01_13_b07e3b4bcffdaab9792bg-066}
\end{center}


\os{newpage} item[{\rot\bf  Digital Data, Analog Signal }]~\\[-3ex]

\bi
  \ite Public telephone system o $300 \mathrm{~Hz}$ to $3400 \mathrm{~Hz}$
\ei
\bi
 \ite Use modem (modulator-demodulator)
 \ei


\bi
  \ite Amplitude shift keying (ASK)

  \ite Frequency shift keying (FSK)

  \ite Phase shift keying (PK)

\ei


\os{newpage} item[{\rot\bf  Modulation Techniques }]~\\[-3ex]

(a) ASK

\begin{center}
\includegraphics[max width=\textwidth]{2023_01_13_b07e3b4bcffdaab9792bg-068}
\end{center}

(b) BFSK

\begin{center}
\includegraphics[max width=\textwidth]{2023_01_13_b07e3b4bcffdaab9792bg-068(1)}
\end{center}

(c) BPSK

\begin{center}
\includegraphics[max width=\textwidth]{2023_01_13_b07e3b4bcffdaab9792bg-068(2)}
\end{center}


\os{newpage} item[{\rot\bf  Amplitude Shift Keying }]~\\[-3ex]

\bi
  \ite Values represented by different amplitudes of carrier

  \ite Usually, one amplitude is zero

\ei
\bi
 \ite i.e. presence and absence of carrier is used
 \ei


\bi
  \ite Susceptible to sudden gain changes

  \ite Inefficient

  \ite Up to 120obps on voice grade lines

  \ite Used over optical fiber

\ei


\os{newpage} item[{\rot\bf  Binary Frequency Shift Keying }]~\\[-3ex]

\bi
  \ite Most common form is binary FSK (BFSK)

  \ite Two binary values represented by two different frequencies (near carrier)

  \ite Less susceptible to error than ASK

  \ite Up to 120obps on voice grade lines

  \ite High frequency radio

  \ite Even higher frequency on LANs using co-ax

\ei


\os{newpage} item[{\rot\bf  Multiple FSK }]~\\[-3ex]

\bi
  \ite More than two frequencies used

  \ite More bandwidth efficient

  \ite More prone to error

  \ite Each signalling element represents more than one bit

\ei


\os{newpage} item[{\rot\bf  FSK on Voice Grade Line }]~\\[-3ex]

signal strength

\begin{center}
\includegraphics[max width=\textwidth]{2023_01_13_b07e3b4bcffdaab9792bg-072}
\end{center}

Figure 5.8 Full-Duplex FSK Transmission on a Voice-Grade Line


\os{newpage} item[{\rot\bf  Phase Shift Keying }]~\\[-3ex]

\bi
  \ite Phase of carrier signal is shifted to represent data - Binary PSK
\ei
\bi
 \ite Two phases represent two binary digits
 \ei


\bi
  \ite Differential PSK
\ei
\bi
 \ite Phase shifted relative to previous transmission rather than some reference signal
 \ei



\os{newpage} item[{\rot\bf  Differential PSK }]~\\[-3ex]

\begin{center}
\includegraphics[max width=\textwidth]{2023_01_13_b07e3b4bcffdaab9792bg-074}
\end{center}


\os{newpage} item[{\rot\bf  Quadrature PSK }]~\\[-3ex]


\os{newpage} item[{\rot\bf  - More efficient use by each signal element representing more than one bit }]~\\[-3ex]
\bi
 \ite e.g. shifts of $\pi / 2\left(90^{\circ}\right)$
 \ei

\bi
 \ite Each element represents two bits
 \ei

\bi
 \ite Can use 8 phase angles and have more than one amplitude o $9600 b p s$ modem use 12 angles, four of which have two amplitudes
 \ei


\bi
  \ite Offset QPSK (orthogonal QPSK)
\ei
\bi
 \ite Delay in Q stream
 \ei



\os{newpage} item[{\rot\bf  QPSK and OQPSK Modulators }]~\\[-3ex]

$\mathrm{I}(\mathrm{t})$

R/2 bps

\begin{center}
\includegraphics[max width=\textwidth]{2023_01_13_b07e3b4bcffdaab9792bg-076}
\end{center}

$a_{n}=\pm 1$

\begin{center}
\includegraphics[max width=\textwidth]{2023_01_13_b07e3b4bcffdaab9792bg-076(1)}
\end{center}

only

\begin{center}
\includegraphics[max width=\textwidth]{2023_01_13_b07e3b4bcffdaab9792bg-076(2)}
\end{center}


\os{newpage} item[{\rot\bf  Examples of QPSF and OQPSK Waveforms }]~\\[-3ex]

\begin{center}
\includegraphics[max width=\textwidth]{2023_01_13_b07e3b4bcffdaab9792bg-077}
\end{center}


\os{newpage} item[{\rot\bf  Performance of Digital to Analog Modulation Schemes }]~\\[-3ex]

\bi
  \ite Bandwidth
\ei

ASK and PSK bandwidth directly related to bit rate

\bi
  \ite FSK bandwidth related to data rate for lower frequencies, but to offset of modulated frequency from carrier at high frequencies
\ei
\bi
 \ite (See Stallings for math)
 \ei


\bi
  \ite In the presence of noise, bit error rate of PSK and QPSK are about $3 \mathrm{~dB}$ superior to $\mathrm{ASK}$ and FSK
\ei


\os{newpage} item[{\rot\bf  Quadrature Amplitude Modulation }]~\\[-3ex]

\bi
  \ite QAM used on asymmetric digital subscriber line (ADSL) and some wireless

  \ite Combination of ASK and PSK

  \ite Logical extension of QPSK

  \ite Send two different signals simultaneously on same carrier frequency

  \ite Use two copies of carrier, one shifted $90^{\circ}$

\ei
\bi
 \ite Each carrier is ASK modulated
 \ei

\bi
 \ite Two independent signals over same medium
 \ei


\bi
  \ite Demodulate and combine for original binary output
\ei


\os{newpage} item[{\rot\bf  QAM Modulator }]~\\[-3ex]

\begin{center}
\includegraphics[max width=\textwidth]{2023_01_13_b07e3b4bcffdaab9792bg-080}
\end{center}

R bps

\begin{center}
\includegraphics[max width=\textwidth]{2023_01_13_b07e3b4bcffdaab9792bg-080(1)}
\end{center}

$\mathrm{d}_{1}(\mathrm{t})$

R/2 bps

\begin{center}
\includegraphics[max width=\textwidth]{2023_01_13_b07e3b4bcffdaab9792bg-080(2)}
\end{center}

$\int \cos 2 \mathrm{pf}_{\mathrm{c}} \mathrm{t}$
\includegraphics[max width=\textwidth, center]{2023_01_13_b07e3b4bcffdaab9792bg-080(3)}

QAM

signal out

$\begin{array}{r}d_{2}(t) \\ R / 2 \mathrm{bps} \\ \hline\end{array}$

$\begin{array}{r}\mathrm{d}_{2}(\mathrm{t}) \\ \mathrm{R} / 2 \mathrm{bps} \\ \hline\end{array}$

\begin{center}
\includegraphics[max width=\textwidth]{2023_01_13_b07e3b4bcffdaab9792bg-080(4)}
\end{center}


\os{newpage} item[{\rot\bf  QAM Levels }]~\\[-3ex]

\bi
  \ite Two level ASK
\ei
\bi
 \ite Each of two streams in one of two states
 \ei

\bi
 \ite Four state system
 \ei

\bi
 \ite Essentially QPSK
 \ei


\bi
  \ite Four level ASK
\ei
\bi
 \ite Combined stream in one of 16 states
 \ei


\bi
  \ite 64 and 256 state systems have been implemented

  \ite Improved data rate for given bandwidth

\ei
\bi
 \ite Increased potential error rate
 \ei



\os{newpage} item[{\rot\bf  Analog Data, Digital Signal }]~\\[-3ex]

\bi
  \ite Digitization

  \ite Conversion of analog data into digital data

  \ite Digital data can then be transmitted using NRZ-L

  \ite Digital data can then be transmitted using code other than NRZ-L

  \ite Digital data can then be converted to analog signal

  \ite Analog to digital conversion done using a codec

  \ite Pulse code modulation

\ei
\bi
 \ite Delta modulation
 \ei



\os{newpage} item[{\rot\bf  Digitizing Analog Data }]~\\[-3ex]

\begin{center}
\includegraphics[max width=\textwidth]{2023_01_13_b07e3b4bcffdaab9792bg-083}
\end{center}

Analog data (voice)
\includegraphics[max width=\textwidth, center]{2023_01_13_b07e3b4bcffdaab9792bg-083(1)}

Digital data

\begin{center}
\includegraphics[max width=\textwidth]{2023_01_13_b07e3b4bcffdaab9792bg-083(2)}
\end{center}

Analog data

(ASK)


\os{newpage} item[{\rot\bf  Pulse Code Modulation(PCM) (1) }]~\\[-3ex]

\bi
  \ite If a signal is sampled at regular intervals at a rate higher than twice the highest signal frequency, the samples contain all the information of the original signal
\ei
\bi
 \ite (Proof - Stallings appendix 4A)
 \ei


\bi
  \ite Voice data limited to below $4000 \mathrm{~Hz}$

  \ite Require 80oo sample per second

  \ite Analog samples (Pulse Amplitude Modulation, PAM)

  \ite Each sample assigned digital value

\ei


\os{newpage} item[{\rot\bf  Pulse Code Modulation(PCM) (2) }]~\\[-3ex]

\bi
  \ite 4 bit system gives 16 levels

  \ite Quantized

\ei
\bi
 \ite Quantizing error or noise
 \ei

\bi
 \ite Approximations mean it is impossible to recover original exactly
 \ei


\bi
  \ite 8 bit sample gives 256 levels

  \ite Quality comparable with analog transmission

  \ite 8000 samples per second of 8 bits each gives $64 \mathrm{kbps}$

\ei


\os{newpage} item[{\rot\bf  PCM Example }]~\\[-3ex]

\begin{center}
\includegraphics[max width=\textwidth]{2023_01_13_b07e3b4bcffdaab9792bg-086}
\end{center}

$\begin{array}{rccccccc}\text { PAM value } & 1.1 & 9.2 & 15.2 & 10.8 & 5.6 & 2.8 & 2.7 \\ \text { quantized code number } & 1 & 9 & 15 & 10 & 5 & 2 & 2 \\ \text { PCM code } & 0001 & 1001 & 1111 & 1010 & 0101 & 0010 & 0010\end{array}$


\os{newpage} item[{\rot\bf  PCM Block Diagram }]~\\[-3ex]

Continuous-time, continuous amplitud (analog) input signal

\begin{center}
\includegraphics[max width=\textwidth]{2023_01_13_b07e3b4bcffdaab9792bg-087}
\end{center}

amplitude signal

(PAM pulses)

\begin{center}
\includegraphics[max width=\textwidth]{2023_01_13_b07e3b4bcffdaab9792bg-087(1)}
\end{center}

(PCM pulses)


\os{newpage} item[{\rot\bf  Nonlinear Encoding }]~\\[-3ex]

\bi
  \ite Quantization levels not evenly spaced

  \ite Reduces overall signal distortion

  \ite Can also be done by companding

\ei


\os{newpage} item[{\rot\bf  Effect of Non-Linear Coding }]~\\[-3ex]

\includegraphics[max width=\textwidth, center]{2023_01_13_b07e3b4bcffdaab9792bg-089}
(a) Without nonlinear encoding
(b) With nonlinear encoding


\os{newpage} item[{\rot\bf  Typical Companding Functions }]~\\[-3ex]

\begin{center}
\includegraphics[max width=\textwidth]{2023_01_13_b07e3b4bcffdaab9792bg-090}
\end{center}


\os{newpage} item[{\rot\bf  Delta Modulation }]~\\[-3ex]

\bi
  \ite Analog input is approximated by a staircase function

  \ite Move up or down one level ( $\delta$ ) at each sample interval

  \ite Binary behavior

  \ite Function moves up or down at each sample interval

\ei


\os{newpage} item[{\rot\bf  Delta Modulation - example }]~\\[-3ex]

\begin{center}
\includegraphics[max width=\textwidth]{2023_01_13_b07e3b4bcffdaab9792bg-092}
\end{center}


\os{newpage} item[{\rot\bf  Delta Modulation - Operation }]~\\[-3ex]

\begin{center}
\includegraphics[max width=\textwidth]{2023_01_13_b07e3b4bcffdaab9792bg-093}
\end{center}

(a) Transmission

\begin{center}
\includegraphics[max width=\textwidth]{2023_01_13_b07e3b4bcffdaab9792bg-093(1)}
\end{center}

(b) Reception


\os{newpage} item[{\rot\bf  Delta Modulation - Performance }]~\\[-3ex]

\bi
  \ite Good voice reproduction

  \ite PCM - 128 levels (7 bit)

\ei
\bi
 \ite Voice bandwidth $4 \mathrm{khz}$
 \ei

\bi
 \ite Should be $8000 \times 7=56 \mathrm{kbps}$ for PCM
 \ei


\bi
  \ite Data compression can improve on this
\ei
\bi
 \ite e.g. Interframe coding techniques for video
 \ei



\os{newpage} item[{\rot\bf  Analog Data, Analog Signals }]~\\[-3ex]

\bi
  \ite Why modulate analog signals?

  \ite Higher frequency can give more efficient transmission

  \ite Permits frequency division multiplexing (chapter 8)

  \ite Types of modulation

  \ite Amplitude

  \ite Frequency

\ei
\bi
 \ite Phase
 \ei



\os{newpage} item[{\rot\bf  Analog Modulation }]~\\[-3ex]

\begin{center}
\includegraphics[max width=\textwidth]{2023_01_13_b07e3b4bcffdaab9792bg-096}
\end{center}

Modulating sine-wave signal

\begin{center}
\includegraphics[max width=\textwidth]{2023_01_13_b07e3b4bcffdaab9792bg-096(1)}
\end{center}

Amplitude-modulated (DSBTC) wave

\begin{center}
\includegraphics[max width=\textwidth]{2023_01_13_b07e3b4bcffdaab9792bg-096(2)}
\end{center}

Phase-modulated wave

\begin{center}
\includegraphics[max width=\textwidth]{2023_01_13_b07e3b4bcffdaab9792bg-096(3)}
\end{center}


\os{newpage} item[{\rot\bf  Ethernet }]~\\[-3ex]


\os{newpage} item[{\rot\bf  Ethernet }]~\\[-3ex]

\bi
  \ite Most successful local area networking technology of last 20 years.

  \ite Developed in the mid-1970s by researchers at the Xerox Palo Alto Research Centers (PARC).

  \ite Uses CSMA/CD technology

  \ite Carrier Sense Multiple Access with Collision Detection.

  \ite A set of nodes send and receive frames over a shared link.

  \ite Carrier sense means that all nodes can distinguish between an idle and a busy link.

  \ite Collision detection means that a node listens as it transmits and can therefore detect when a frame it is transmitting has collided with a frame transmitted by another node.

\ei


\os{newpage} item[{\rot\bf  Ethernet }]~\\[-3ex]

\bi
  \ite Uses ALOHA (packet radio network) as the root protocol

  \ite Developed at the University of Hawaii to support communication across the Hawaiian Islands.

  \ite For ALOHA the medium was atmosphere, for Ethernet the medium is a coax cable.

  \ite DEC and Intel joined Xerox to define a 10-Mbps Ethernet standard in 1978.

  \ite This standard formed the basis for IEEE standard 802.3

  \ite More recently 802.3 has been extended to include a 100Mbps version called Fast Ethernet and a 1000-Mbps version called Gigabit Ethernet.

\ei


\os{newpage} item[{\rot\bf  Ethernet }]~\\[-3ex]

\bi
  \ite An Ethernet segment is implemented on a coaxial cable of up to $500 \mathrm{~m}$.

  \ite This cable is similar to the type used for cable TV except that it typically has an impedance of 50 ohms instead of cable TV's 75 ohms.

  \ite Hosts connect to an Ethernet segment by tapping into it.

  \ite A transceiver (a small device directly attached to the tap) detects when the line is idle and drives signal when the host is transmitting.

  \ite The transceiver also receives incoming signal.

  \ite The transceiver is connected to an Ethernet adaptor which is plugged into the host.

  \ite The protocol is implemented on the adaptor.

\ei


\os{newpage} item[{\rot\bf  Ethernet }]~\\[-3ex]

Transceiver

\begin{center}
\includegraphics[max width=\textwidth]{2023_01_13_b07e3b4bcffdaab9792bg-101}
\end{center}

Adaptor

Host


\os{newpage} item[{\rot\bf  Ethernet transceiver and adaptor }]~\\[-3ex]


\os{newpage} item[{\rot\bf  Ethernet }]~\\[-3ex]

\bi
  \ite Multiple Ethernet segments can be joined together by repeaters.

  \ite A repeater is a device that forwards digital signals.

  \ite No more than four repeaters may be positioned between any pair of hosts.

\ei
\bi
 \ite An Ethernet has a total reach of only $2500 \mathrm{~m}$.
 \ei



\os{newpage} item[{\rot\bf  Ethernet - Repeater }]~\\[-3ex]

\begin{center}
\includegraphics[max width=\textwidth]{2023_01_13_b07e3b4bcffdaab9792bg-103}
\end{center}

\begin{center}
\includegraphics[max width=\textwidth]{2023_01_13_b07e3b4bcffdaab9792bg-103(1)}
\end{center}

Repeater

\begin{center}
\includegraphics[max width=\textwidth]{2023_01_13_b07e3b4bcffdaab9792bg-103(2)}
\end{center}


\os{newpage} item[{\rot\bf  Ethernet }]~\\[-3ex]

\bi
  \ite Any signal placed on the Ethernet by a host is broadcast over the entire network
\ei
\bi
 \ite Signal is propagated in both directions.
 \ei

\bi
 \ite Repeaters forward the signal on all outgoing segments.
 \ei

\bi
 \ite Terminators attached to the end of each segment absorb the signal.
 \ei


\bi
  \ite Ethernet uses Manchester encoding scheme.
\ei


\os{newpage} item[{\rot\bf  Ethernet }]~\\[-3ex]

\bi
  \ite New Technologies in Ethernet

  \ite Instead of using coax cable, an Ethernet can be constructed from a thinner cable known as 10 Base2 (the original was 10Base5)

\ei

10 means the network operates at $10 \mathrm{Mbps}$

· Base means the cable is used in a baseband system

$\times 2$ means that a given segment can be no longer than $200 \mathrm{~m}$


\os{newpage} item[{\rot\bf  Ethernet }]~\\[-3ex]

\bi
  \ite New Technologies in Ethernet - Another cable technology is 10BaseT
\ei

\begin{CJK}{UTF8}{mj}乂\end{CJK} T stands for twisted pair

\bi
  \ite Limited to $100 \mathrm{~m}$ in length
\ei
\bi
 \ite With 10 BaseT, the common configuration is to have several point to point segments coming out of a multiway repeater, called $H u b$
 \ei



\os{newpage} item[{\rot\bf  Ethernet }]~\\[-3ex]

\begin{center}
\includegraphics[max width=\textwidth]{2023_01_13_b07e3b4bcffdaab9792bg-107}
\end{center}

Ethernet Hub


\os{newpage} item[{\rot\bf  Access Protocol for Ethernet }]~\\[-3ex]

\bi
  \ite The algorithm is commonly called Ethernet's Media Access Control (MAC).
\ei
\bi
 \ite It is implemented in Hardware on the network adaptor.
 \ei


\bi
  \ite Frame format
\ei
\bi
 \ite Preamble (64bit): allows the receiver to synchronize with the signal (sequence of alternating os and 1s).
 \ei

\bi
 \ite Host and Destination Address (48bit each).
 \ei

\bi
 \ite Packet type (16bit): acts as demux key to identify the higher level protocol.
 \ei

\bi
 \ite Data (up to 1500 bytes)
 \ei


\bi
  \ite Minimally a frame must contain at least 46 bytes of data.

  \ite Frame must be long enough to detect collision.

\ei
\bi
 \ite CRC (32bit)
 \ei



\os{newpage} item[{\rot\bf  Ethernet Frame }]~\\[-3ex]

\begin{center}
\begin{tabular}{|c|c|c|c|c|c|}
64 & 48 & 48 & 16 & 32 &  \\
\hline
Preamble & Dest addr & Src addr & Type & Body & CRC \\
\hline
\end{tabular}
\end{center}


\os{newpage} item[{\rot\bf  Ethernet Frame Format }]~\\[-3ex]


\os{newpage} item[{\rot\bf  Ethernet Addresses }]~\\[-3ex]

\bi
  \ite Each host on an Ethernet (in fact, every Ethernet host in the world) has a unique Ethernet Address.

  \ite The address belongs to the adaptor, not the host.

\ei
\bi
 \ite It is usually burnt into ROM.
 \ei


\bi
  \ite Ethernet addresses are typically printed in a human readable format
\ei
\bi
 \ite As a sequence of six numbers separated by colons.
 \ei

\bi
 \ite Each number corresponds to 1 byte of the 6 byte address and is given by a pair of hexadecimal digits, one for each of the 4-bit nibbles in the byte
 \ei

\bi
 \ite Leading os are dropped.
 \ei

\bi
 \ite For example, 8:0:2b:e4:b1:2 is
 \ei


· 000010000000000000101011111001001011000100000010


\os{newpage} item[{\rot\bf  Ethernet Addresses }]~\\[-3ex]

\bi
  \ite To ensure that every adaptor gets a unique address, each manufacturer of Ethernet devices is allocated a different prefix that must be prepended to the address on every adaptor they build
\ei

\bi
  \ite AMD has been assigned the 24bit prefix 8:0:20
\ei


\os{newpage} item[{\rot\bf  Ethernet Addresses }]~\\[-3ex]

\bi
  \ite Each frame transmitted on an Ethernet is received by every adaptor connected to that Ethernet.

  \ite Each adaptor recognizes those frames addressed to its address and passes only those frames on to the host.

  \ite In addition, to unicast address, an Ethernet address consisting of all $1 \mathrm{~s}$ is treated as a broadcast address.

\ei
\bi
 \ite All adaptors pass frames addressed to the broadcast address up to the host.
 \ei


\bi
  \ite Similarly, an address that has the first bit set to 1 but is not the broadcast address is called a multicast address.
\ei
\bi
 \ite A given host can program its adaptor to accept some set of multicast addresses.
 \ei



\os{newpage} item[{\rot\bf  Ethernet Addresses }]~\\[-3ex]

\bi
  \ite To summarize, an Ethernet adaptor receives all frames and accepts
\ei
\bi
 \ite Frames addressed to its own address
 \ei

\bi
 \ite Frames addressed to the broadcast address
 \ei

\bi
 \ite Frames addressed to a multicast addressed if it has been instructed
 \ei



\os{newpage} item[{\rot\bf  Ethernet Transmitter Algorithm }]~\\[-3ex]

\bi
  \ite When the adaptor has a frame to send and the line is idle, it transmits the frame immediately.

  \ite The upper bound of 1500 bytes in the message means that the adaptor can occupy the line for a fixed length of time.

  \ite When the adaptor has a frame to send and the line is busy, it waits for the line to go idle and then transmits immediately.

  \ite The Ethernet is said to be 1-persistent protocol because an adaptor with a frame to send transmits with probability 1 whenever a busy line goes idle.

\ei


\os{newpage} item[{\rot\bf  Ethernet Transmitter Algorithm }]~\\[-3ex]

\bi
  \ite Since there is no centralized control it is possible for two (or more) adaptors to begin transmitting at the same time,

  \ite Either because both found the line to be idle,

  \ite Or, both had been waiting for a busy line to become idle.

  \ite When this happens, the two (or more) frames are said to be collide on the network.

\ei


\os{newpage} item[{\rot\bf  Ethernet Transmitter Algorithm }]~\\[-3ex]

\bi
  \ite Since Ethernet supports collision detection, each sender is able to determine that a collision is in progress.

  \ite At the moment an adaptor detects that its frame is colliding with another, it first makes sure to transmit a 32-bit jamming sequence and then stops transmission.

  \ite Thus, a transmitter will minimally send 96 bits in the case of collision

\ei

64-bit preamble $+32$-bit jamming sequence


\os{newpage} item[{\rot\bf  Ethernet Transmitter Algorithm }]~\\[-3ex]

\bi
  \ite One way that an adaptor will send only 96 bit (called a runt frame) is if the two hosts are close to each other.

  \ite Had they been farther apart,

  \ite They would have had to transmit longer, and thus send more bits, before detecting the collision.

\ei


\os{newpage} item[{\rot\bf  Ethernet Transmitter Algorithm }]~\\[-3ex]

\bi
  \ite The worst case scenario happens when the two hosts are at opposite ends of the Ethernet.

  \ite To know for sure that the frame its just sent did not collide with another frame, the transmitter may need to send as many as 512 bits.

\ei
\bi
 \ite Every Ethernet frame must be at least 512 bits (64 bytes) long.
 \ei


$\times 14$ bytes of header $+46$ bytes of data $+4$ bytes of CRC


\os{newpage} item[{\rot\bf  Ethernet Transmitter Algorithm }]~\\[-3ex]

\bi
  \ite Why 512 bits?
\ei
\bi
 \ite Why is its length limited to $2500 \mathrm{~m}$ ?
 \ei



\os{newpage} item[{\rot\bf  - The farther apart two nodes are, the longer it takes for a frame sent by one to reach the other, and the network is vulnerable to collision during this time }]~\\[-3ex]


\os{newpage} item[{\rot\bf  Ethernet Transmitter Algorithm }]~\\[-3ex]

\bi
  \ite A begins transmitting a frame at time $t$

  \ite $d$ denotes the one link latency

  \ite The first bit of A's frame arrives at B at time $t+d$

  \ite Suppose an instant before host A's frame arrives, host B begins to transmit its own frame

  \ite B's frame will immediately collide with A's frame and this collision will be detected by host B

  \ite Host B will send the 32-bit jamming sequence

  \ite Host A will not know that the collision occurred until B's frame reaches it, which will happen at $t+2 * d$

  \ite Host A must continue to transmit until this time in order to detect the collision

  \ite Host A must transmit for $2{ }^{*} d$ to be sure that it detects all possible collisions

\ei


\os{newpage} item[{\rot\bf  Ethernet Transmitter Algorithm }]~\\[-3ex]

(a)

\begin{center}
\includegraphics[max width=\textwidth]{2023_01_13_b07e3b4bcffdaab9792bg-121}
\end{center}

(b)

\begin{center}
\includegraphics[max width=\textwidth]{2023_01_13_b07e3b4bcffdaab9792bg-121(1)}
\end{center}

(c)

\begin{center}
\includegraphics[max width=\textwidth]{2023_01_13_b07e3b4bcffdaab9792bg-121(2)}
\end{center}

(d) $\mathrm{A}$

\begin{center}
\includegraphics[max width=\textwidth]{2023_01_13_b07e3b4bcffdaab9792bg-121(3)}
\end{center}

Worst-case scenario:

(a) A sends a frame at time $t$;

(b) A's frame arrives at $\mathrm{B}$ at time $t+d$;

(c) B begins transmitting at time $t+d$ and collides with A's frame;

(d) B's runt (32-bit) frame arrives at $\mathrm{A}$ at time $t+2 d$.


\os{newpage} item[{\rot\bf  Ethernet Transmitter Algorithm }]~\\[-3ex]


\os{newpage} item[{\rot\bf  - Consider that a maximally configured Ethernet is 2500 $\mathrm{m}$ long, and there may be up to four repeaters between any two hosts, the round trip delay has been determined to be $51.2 \mu \mathrm{s }]~\\[-3ex]
$
 o Which on $10 \mathrm{Mbps}$ Ethernet corresponds to 512 bits}

\os{newpage} item[{\rot\bf  - The other way to look at this situation, }]~\\[-3ex]
\bi
 \ite We need to limit the Ethernet's maximum latency to a fairly small value $(51.2 \mu \mathrm{s})$ for the access algorithm to work
 \ei


\bi
  \ite Hence the maximum length for the Ethernet is on the order of $\mathbf{2 5 0 0}$ $\mathrm{m}$.
\ei


\os{newpage} item[{\rot\bf  Ethernet Transmitter Algorithm }]~\\[-3ex]

\bi
  \ite Once an adaptor has detected a collision, and stopped its transmission, it waits a certain amount of time and tries again.

  \ite Each time the adaptor tries to transmit but fails, it doubles the amount of time it waits before trying again.

  \ite This strategy of doubling the delay interval between each retransmission attempt is known as Exponential Backoff.

\ei


\os{newpage} item[{\rot\bf  Ethernet Transmitter Algorithm }]~\\[-3ex]

\bi
  \ite The adaptor first delays either o or $51.2 \mu \mathrm{s}$, selected at random.

  \ite If this effort fails, it then waits $0,51.2,102.4,153.6 \mu \mathrm{s}$ (selected randomly) before trying again;

  \ite This is $k^{*} 51.2$ for $k=0,1,2,3$

  \ite After the third collision, it waits $k^{*} 51.2$ for $k=0 \ldots 2^{3}-1$ (again selected at random).

  \ite In general, the algorithm randomly selects a $k$ between $o$ and $2^{\mathrm{n}}-1$ and waits for $k^{*} 51.2 \mu \mathrm{s}$, where $n$ is the number of collisions experienced so far.

\ei


\os{newpage} item[{\rot\bf  Experience with Ethernet }]~\\[-3ex]

\bi
  \ite Ethernets work best under lightly loaded conditions.

  \ite Under heavy loads, too much of the network's capacity is wasted by collisions.

  \ite Most Ethernets are used in a conservative way.

  \ite Have fewer than 200 hosts connected to them which is far fewer than the maximum of 1024.

  \ite Most Ethernets are far shorter than $2500 m$ with a round-trip delay of closer to $5 \mu$ s than $51.2 \mu \mathrm{s}$.

  \ite Ethernets are easy to administer and maintain.

  \ite There are no switches that can fail and no routing and configuration tables that have to be kept up-to-date.

  \ite It is easy to add a new host to the network.

  \ite It is inexpensive.

\ei

C Cable is cheap, and only other cost is the network adaptor on each host.


\os{newpage} item[{\rot\bf  Orientation }]~\\[-3ex]

1

2

\bi
  \ite IP (Internet Protocol) is a Network Layer Protocol.
\ei

\begin{center}
\includegraphics[max width=\textwidth]{2023_01_13_b07e3b4bcffdaab9792bg-127}
\end{center}

\bi
  \ite IP's current version is Version 4 (IPv4). It is specified in RFC 891.
\ei


\os{newpage} item[{\rot\bf  IP: The waist of the hourglass }]~\\[-3ex]

1

2

\bi
  \ite IP is the waist of the liourglass of the Internet protocol architecture

  \ite Multiple higher-layer protocols

  \ite Multiple lower-layer protocols

  \ite Only one protocol at the network

\ei

\begin{center}
\includegraphics[max width=\textwidth]{2023_01_13_b07e3b4bcffdaab9792bg-128}
\end{center}


\os{newpage} item[{\rot\bf  Application protocol }]~\\[-3ex]

1

2

\bi
  \ite IP is the highest layer protocol which is implemented at both routers and hosts
\ei

\begin{center}
\includegraphics[max width=\textwidth]{2023_01_13_b07e3b4bcffdaab9792bg-129}
\end{center}


\os{newpage} item[{\rot\bf  IP Service }]~\\[-3ex]

\bi
  \ite Delivery service of IP is minimal 1
3
0

  \ite IP provide provides an unreliable connectionless best effort service (also called: "datagram service").

\ei
\bi
 \ite Unreliable: IP does not make an attempt to recover lost packets
 \ei

\bi
 \ite Connectionless: Each packet ("datagram") is handled independently. IP is not aware that packets between hosts may be sent in a logical sequence
 \ei

\bi
 \ite Best effort: IP does not make guarantees on the service (no throughput guarantee, no delay guarantee,..)
 \ei



\os{newpage} item[{\rot\bf  - Consequences: }]~\\[-3ex]

\bi
  \ite Higher layer protocols have to deal with losses or with duplicate packets

  \ite Packets may be delivered out-of-sequence

\ei


\os{newpage} item[{\rot\bf  IP Service }]~\\[-3ex]

1

3

\bi
  \ite IP supports the following services:
\ei

ane-to-one

cone-to-all

rone-to-several

\begin{center}
\includegraphics[max width=\textwidth]{2023_01_13_b07e3b4bcffdaab9792bg-131}
\end{center}

unicast

\begin{center}
\includegraphics[max width=\textwidth]{2023_01_13_b07e3b4bcffdaab9792bg-131(1)}
\end{center}

broadcast (unicast)

(broadcast)

(multicast)

\begin{center}
\includegraphics[max width=\textwidth]{2023_01_13_b07e3b4bcffdaab9792bg-131(2)}
\end{center}

\begin{center}
\includegraphics[max width=\textwidth]{2023_01_13_b07e3b4bcffdaab9792bg-131(3)}
\end{center}

\bi
  \ite IP multicast also supports a many-to-many service.

  \ite IP multicast requires support of other protocols (IGMP, multicast routing)

\ei

\os{newpage} item[{\rot\bf IP Datagram Format }]~\\[-3ex]


\bi
  \ite 20 bytes $\leq$ Header Size $<2^{4}$ x 4 bytes $=60$ bytes

  \ite 20 bytes $\leq$ Total Length $<2^{16}$ bytes $=65536$ bytes

\ei

\begin{center}
\includegraphics[max width=\textwidth]{2023_01_13_b07e3b4bcffdaab9792bg-132}
\end{center}


\os{newpage} item[{\rot\bf  IP Datagram Format }]~\\[-3ex]

$$
1
$$

3

\bi
  \ite Question: In which order are the bytes of an IP datagram transmitted?

  \ite Answer:

  \ite Transmission is row by row

  \ite For each row:

\ei

\begin{enumerate}
  \ite First transmit bits $0-7$

  \ite Then transmit bits $8-15$

  \ite Then transmit bits 16-23

  \ite Then transmit bits 24-31

\end{enumerate}


\os{newpage} item[{\rot\bf  - This is called network byte order or big endian byte ordering. }]~\\[-3ex]

\bi
  \ite Note: Many computers (incl. Intel processors) store 32-bit words in little endian format. Others (incl. Motorola processors) use big endian.
\ei


\os{newpage} item[{\rot\bf  Big endian vs. small endian }]~\\[-3ex]


\bi
  \ite Conventions to store a multibyte work

  \ite Example: a 4 byte Long Integer Byte3 Byte2 Byte1 Byte0

\ei


\os{newpage} item[{\rot\bf  Little Endian }]~\\[-3ex]

\bi
  \ite Stores the low-order byte at the lowest address and the highest order byte in the highest address.
\ei

Base Addresst0 Byte0 Base Addresst1 Byte1 Base Address+2 Byte2 Base Address+3 Byte3

\bi
  \ite Intel processors use this order
\ei


\os{newpage} item[{\rot\bf  Big Endian }]~\\[-3ex]

\bi
  \ite Stores the high-order byte at the lowest address, and the low-order byte at the highest address.
Base Address+0 Byte3
Base Addresst1 Byte2
Base Addresst2 Byte1
Base Address+3 Byte0
\ei

Motorola processors use big endian.

\os{newpage} item[{\rot\bf Fields of the IP Header } ]~\\[-3ex]


\bi
  \ite Version (4 bits): current version is 4, next version will be 6.

  \ite Header length (4 bits): length of IP header, in multiples of 4 bytes

  \ite DS/ECN field (1 byte)

\ei
\bi
 \ite This field was previously called as Type-of-Service (TOS) field. The role of this field has been re-defined, but is "backwards compatible" to TOS interpretation
 \ei

\bi
 \ite Differentiated Service (DS) (6 bits):
 \ei


\bi
  \ite Used to specify service level (currently not supported in the Internet)
\ei
\bi
 \ite Explicit Congestion Notification (ECN) (2 bits):
 \ei


\begin{CJK}{UTF8}{mj}乂\end{CJK} New feedback mechanism used by TCP

\os{newpage} item[{\rot\bf Fields of the IP Header } ]~\\[-3ex]


\bi
  \ite Identification (16 bits): Unique identification of a datagram from a host. Incremented whenever a datagram is transmitted

  \ite Flags (3 bits):

  \ite First bit always set to 0

\ei
\bi
 \ite DF bit (Do not fragment)
 \ei

\bi
 \ite MF bit (More fragments)
 \ei


Will be explained later $\rightarrow$ Fragmentation


\os{newpage} item[{\rot\bf  Fields of the IP Header }]~\\[-3ex]

1

3

\bi
  \ite Time To Live (TTL) (17byte):
\ei
\bi
 \ite Specifies longest paths before datagram is dropped
 \ei

\bi
 \ite Role of TTL field: Ensure that packet is eventually dropped when a routing loop occurs
 \ei


Used as follows:
\bi
 \ite Sender sets the value (e.g., 64)
 \ei

\bi
 \ite Each router decrements the value by 1
 \ei

\bi
 \ite When the value reaches 0 , the datagram is dropped
 \ei



\os{newpage} item[{\rot\bf  Fields of the IP Header }]~\\[-3ex]

1

\bi
  \ite Protocol (1 byte):
\ei

3

8

\bi
  \ite Specifies the higher-layer protocol.
\ei

\bi
  \ite Used for demultiplexing to higher layers.
\ei

\begin{center}
\includegraphics[max width=\textwidth]{2023_01_13_b07e3b4bcffdaab9792bg-138}
\end{center}

\bi
  \ite Header checksum (2 bytes): A simple 16-bit long checksum which is computed for the header of the datagram.
\ei

\os{newpage} item[{\rot\bf   Fields of the IP Header } ]~\\[-3ex]



\os{newpage} item[{\rot\bf  - Options: }]~\\[-3ex]

9

\bi
  \ite Security restrictions
\ei

Record Route: each router that processes the packet adds its IP address to the header.

Timestamp: each router that processes the packet adds its IP address and time to the header.

(loose) Source Routing: specifies a list of routers that must be traversed.

\bi
  \ite (strict) Source Routing: specifies a list of the only routers that can be traversed.
\ei


\os{newpage} item[{\rot\bf  - Padding: Padding bytes are added to ensure that header ends on a 4-byte boundary }]~\\[-3ex]


\os{newpage} item[{\rot\bf  Maximum Transmission Unit }]~\\[-3ex]

1

4

\bi
  \ite Maximum size of IP datagram is 65535 , but the data link layer protocol generally imposes a limit that is much smaller

  \ite Example:

\ei
\bi
 \ite Ethernet frames have a maximum payload of 1500 bytes
 \ei


$\rightarrow$ IP datagrams encapsulated in Ethernet frame cannot be longer than 1500 bytes

\bi
  \ite The limit on the maximum IP datagram size, imposed by the data link protocol is called maximum transmission unit (MTU)

  \ite MTUs for various data link protocols:

\ei

Ethernet: $\quad 1500$

802.3:

802.5: 1492

4464 FDDI: $\quad 4352$

ATM AAL5: 9180

PPP:


\item[{\rot\bf  IP Fragmentation  }]~\\[-3ex]
 ]~\\[-3ex]


\bi
  \ite What if the size of an IP datagram exceeds the MTU? IP datagram is fragmented into smaller units.

  \ite What if the route contains networks with different MTUs?

\ei

\begin{center}
\includegraphics[max width=\textwidth]{2023_01_13_b07e3b4bcffdaab9792bg-141}
\end{center}

MTUs: FDDI: 4352

Ethernet: 1500

\bi
  \ite Fragmentation:

  \ite IP router splits the datagram into several datagram

  \ite Fragments are reassembled at receiver

\ei


\item[{\rot\bf  Where is Fragmentation done?  }]~\\[-3ex]



\bi
  \ite Fragmentation can be done at the sender or at intermediate routers
\ei

×he same datagram can be fragmented several times.

 Reassembly of original datagram is only done at destination hosts !!
\includegraphics[max width=\textwidth, center]{2023_01_13_b07e3b4bcffdaab9792bg-142}

Fragment 2

42

Fragment 1

(1)

Router


\os{newpage} item[{\rot\bf  What's involved in Fragmentation? }]~\\[-3ex]

\bi
  \ite The following fields in the IP header are involved:
\ei

\begin{center}
\includegraphics[max width=\textwidth]{2023_01_13_b07e3b4bcffdaab9792bg-143}
\end{center}

Identification

Flags When a datagram is fragmented, the identification is the same in all fragments

DF bit is set: Datagram cannot be fragmented and must be discarded if MTU is too small

MF bit set: This datagram is part of a fragment and an additional fragment follows this one


\os{newpage} item[{\rot\bf  What's involved in Fragmentation? }]~\\[-3ex]

\bi
  \ite The following fields in the IP header are involved:
\ei

\begin{center}
\includegraphics[max width=\textwidth]{2023_01_13_b07e3b4bcffdaab9792bg-144}
\end{center}

Fragment offset Offset of the payload of the current fragment in the original datagram Total length

Total length of the current fragment


\item[{\rot\bf  Example of Fragmentation  }]~\\[-3ex]



\bi
  \ite A datagram with size 2400 bytes isust be fragmented according to an MTU limit of 1000 bytes
\ei

Header length: 20 Total length: 2400 Identification: 0xa428

DF flag: 0 MF flag: 0

Fragment offset: 0 Header length: 20

Total length: 448 Identification: 0xa428

DF flag: 0

MF flag: 0

Fragment offset: 244 Header length: 20

Total length: 996 Identification: 0xa428 DF flag: 0 MF flag: 1 Fragment offset: 122 Header length: 20 Total length: 996 Identification: 0xa428 DF flag: 0 MF flag: 1 fragment offset: 0
\includegraphics[max width=\textwidth, center]{2023_01_13_b07e3b4bcffdaab9792bg-145}

Fragment 2

Fragment 1

Router


\os{newpage} item[{\rot\bf  Determining the length of fragments }]~\\[-3ex]

1

4

\bi
  \ite To determine the size of the fragments we recall that, since there are only 13 bits available for the fragment offset, the offset is given as a multiple of eight bytes.

  \ite a result, the first and second fragment have a size of 996 bytes (and not 1000 bytes). This number is chosen since 976 is the largest number smaller than $1000-20=980$ that is divisible by eight.

  \ite The payload for the first and second fragments is 976 bytes long, with bytes o through 975 of the original IP payload in the first fragment, and bytes 976 through 1951 in the second fragment.

  \ite The payload of the third fragment has the remaining 428 bytes, from byte 1952 through 2379.

  \ite With these considerations, we can determine the values of the fragment offset, which are $0,976 / 8=122$, and $1952 / 8=$ 244 , respectively, for the first, second and third fragment.

\ei


\os{newpage} item[{\rot\bf  IP Addressing }]~\\[-3ex]


\item[{\rot\bf  What is an IP address?  }]~\\[-3ex]


\os{newpage} item[{\rot\bf  - An unique identifier for a computer or device (host) on a TCP/IP network }]~\\[-3ex]


\os{newpage} item[{\rot\bf  - A 32-bit binary number usually represented as 4 decimal numbers separated by a period }]~\\[-3ex]




\item[{\rot\bf  What is an IP address?  }]~\\[-3ex]


\os{newpage} item[{\rot\bf  - Each address is $3^{2}$ bits wide }]~\\[-3ex]

\bi
  \ite Valid addresses can range from 0.0.0.0 to 255.255.255.255
\ei




\item[{\rot\bf  What is an IP address?  }]~\\[-3ex]


\os{newpage} item[{\rot\bf  Theoretically, a total of $\approx 4.3$ billion addresses are available }]~\\[-3ex]




\os{newpage} item[{\rot\bf  Two addresses in one... }]~\\[-3ex]


\os{newpage} item[{\rot\bf  Each address consists of two parts }]~\\[-3ex]


\os{newpage} item[{\rot\bf  The network address }]~\\[-3ex]


\os{newpage} item[{\rot\bf  The host address }]~\\[-3ex]


\os{newpage} item[{\rot\bf  Other systems may use more than one address (Ex: IPX) }]~\\[-3ex]


\os{newpage} item[{\rot\bf  The Five Network Classes }]~\\[-3ex]

\begin{enumerate}
  \ite Class $\mathrm{A}-$ begins with 0
\end{enumerate}

\bi
  \ite $00000001\left(1_{10}\right)$ to $01111111\left(126_{10}\right)^{*}$
\ei

\begin{enumerate}
  \setcounter{enumi}{1}
  \ite Class $\mathbf{B}-$ begins with $\mathbf{1 0}$
\end{enumerate}

\bi
  \ite $10000000\left(128_{10}\right)$ to $10111111\left(191_{10}\right)$
\ei


\os{newpage} item[{\rot\bf  Class $\mathrm{C}-$ begins with $\mathbf{1 1 0}$ }]~\\[-3ex]

\bi
  \ite $11000000\left(192_{10}\right)$ to $11011111\left(223_{10}\right)$
\ei

$* 01111111=127_{10}$


\os{newpage} item[{\rot\bf  Addresses beginning with 127 are reserved for loopback (127.0.0.1 is YoU) }]~\\[-3ex]


\os{newpage} item[{\rot\bf  The Five Network Classes }]~\\[-3ex]


\os{newpage} item[{\rot\bf  Class $\mathrm{D}$ - begins with 1110 }]~\\[-3ex]

\bi
  \ite $224_{10}$ to $239_{10}$

  \ite Reserved for multicasting

\ei


\os{newpage} item[{\rot\bf  Class E - begins with $\mathbf{1 1 1 1}$ }]~\\[-3ex]

\bi
  \ite $240_{10}$ to $254_{10}$

  \ite Reserved for future use

\ei


\os{newpage} item[{\rot\bf  These should not be used for host addressing }]~\\[-3ex]


\item[{\rot\bf  Which part belongs to the network and which part belongs to the node?  }]~\\[-3ex]



\os{newpage} item[{\rot\bf  Class B - XXXXXXXX.XXXXXXXX.yyyyyyyy.yyyyyyyy }]~\\[-3ex]


\os{newpage} item[{\rot\bf  Class C - XXXXXXXX.XXXXXXXX.XXXXXXXX.yyyyyyyy }]~\\[-3ex]


\os{newpage} item[{\rot\bf  Where $X=$ Network and $y=$ node }]~\\[-3ex]


\os{newpage} item[{\rot\bf  IP Addresses* }]~\\[-3ex]


\os{newpage} item[{\rot\bf  Class $1^{\text {st }}$ Octet Networks Ids Host IDs }]~\\[-3ex]

\begin{center}
\begin{tabular}{|c|c|c|c|}
\hline
$A$ & $1-126$ & $2^{7}=126$ & $2^{24}=16 M$ \\
\hline
$B$ & $128-191$ & $2^{14}=16 K$ & $2^{16}=64 K$ \\
\hline
$C$ & $192-223$ & $2^{21}=2 M$ & $2^{8}=255$ \\
\hline
\end{tabular}
\end{center}


\os{newpage} item[{\rot\bf  There are three IP network addresses reserved for private networks }]~\\[-3ex]

\begin{enumerate}
  \ite $10.0 .0 .0 / 8$

  \ite $172.16 .0 .0 / 12$

  \ite $192.168 .0 .0 / 16$

\end{enumerate}


\os{newpage} item[{\rot\bf  These can be used by anyone setting up an internal network. }]~\\[-3ex]


\os{newpage} item[{\rot\bf  Routers will never forward packets coming from these addresses. }]~\\[-3ex]


\os{newpage} item[{\rot\bf  Subnetting }]~\\[-3ex]

...can be done for a variety of reasons
\bi
 \ite Organization
 \ei


\bi
  \ite Use of different physical media

  \ite Preservation of address space

\ei
\bi
 \ite Security
 \ei


\bi
  \ite The most common reason is to control network traffic
\ei


\os{newpage} item[{\rot\bf  Subnetting }]~\\[-3ex]

\bi
  \ite In an Ethernet network, all nodes on a segment see all packets transmitted by other nodes on that segment

  \ite Performance can be adversely affected under heavy traffic loads

  \ite A router is used to connect IP networks to minimize the amount of traffic each segment must receive

\ei


\os{newpage} item[{\rot\bf  Subnet masking }]~\\[-3ex]

\bi
  \ite Applying a subnet mask allows you to identify the network and node parts of the address. A router will then determine whether the address is local or remote.

  \ite Network bits are masked as 1s

  \ite Node bits are masked as os

  \ite Class A - 255.0.0.0

\ei
\bi
 \ite 11111111.00000000.00000000.00000000
 \ei


\bi
  \ite Class B - 255.255.0.0
\ei

O 11111111.11111111.00000000.00000000

\bi
  \ite Class C - 255.255.255.0
\ei
\bi
 \ite 11111111.11111111.11111111.00000000
 \ei



\os{newpage} item[{\rot\bf  Subnet masking }]~\\[-3ex]

cil c: WINDOWS system32 cmd.exe

$P: \backslash>$ ipconfig $/ \mathbf{a l l}$

Windows IP Conf iguration

Host Name : - $-\cdot \cdot \cdot \cdot \cdot \cdot \cdot \cdot \cdot \cdot=$ D201B-01

\begin{center}
\includegraphics[max width=\textwidth]{2023_01_13_b07e3b4bcffdaab9792bg-160}
\end{center}

Node Type ... . . . . . . . . Hybrid

IP Routing Enabled: : : : : : : No

WINS Proxy Enabled.: :.:-:.: No

DNS Suffix Search List... . . : \href{http://Hrazosport.cc.tx.us}{Hrazosport.cc.tx.us}

\href{http://cc.tx.us}{cc.tx.us}

\href{http://tx.us}{tx.us}

Ethernet adapter Local Area Connection:

Connection-specific DNS Suffix

\begin{center}
\includegraphics[max width=\textwidth]{2023_01_13_b07e3b4bcffdaab9792bg-160(1)}
\end{center}

rollers

Physical Address . . . . . . . . . : 00-19-B9-46-38-F2

Dhcp Enabled. . . . . . . . . . . : No

\begin{center}
\includegraphics[max width=\textwidth]{2023_01_13_b07e3b4bcffdaab9792bg-160(2)}
\end{center}

\begin{center}
\includegraphics[max width=\textwidth]{2023_01_13_b07e3b4bcffdaab9792bg-160(3)}
\end{center}

\begin{center}
\includegraphics[max width=\textwidth]{2023_01_13_b07e3b4bcffdaab9792bg-160(4)}
\end{center}

DNS Seruers. $=.-.=.-.=-206.40 .177 .10$

$206.40 .177 .11$

$P: \backslash>$


\os{newpage} item[{\rot\bf  Subnet masking }]~\\[-3ex]

\bi
  \ite Performing a bitwise logical AND between the IP address and the subnet mask results in the network address

  \ite Ex: Class - B 140.179.240.200

\ei

10001100.10110011.11110000.11001000

11111111.11111111.00000000.00000000

10001100.10110011.00000000.00000000

Network Address $=140.179 .000 .000$


\item[{\rot\bf  A Few Rules...  }]~\\[-3ex]


\os{newpage} item[{\rot\bf  Each device on a node has a unique MAC address }]~\\[-3ex]


\os{newpage} item[{\rot\bf  Each device on a node needs a unique IP address }]~\\[-3ex]


\os{newpage} item[{\rot\bf  All devices on the same physical segment share a common network ID (subnet mask) }]~\\[-3ex]




\os{newpage} item[{\rot\bf  Address Resolution Protocol (ARP) }]~\\[-3ex]

\bi
  \ite Before an IP packet can be forwarded to another host, the MAC address (usually 6 bytes written in hex (Ex: 02-FE-874A-8C-A9) of the receiving machine must be known

  \ite ARP determines the MAC addresses that correspond to an IP address

  \ite A router will choose direct paths for the network packets based on the addressing of the IP frame it is handling (different routes to different networks)

\ei


\os{newpage} item[{\rot\bf  Address Resolution Protocol (ARP) }]~\\[-3ex]


\item[{\rot\bf  Internet and Data Link Layer Addresses  }]~\\[-3ex]

\bi
  \ite Each host and router on a subnet needs a data link layer address to specify its address on the subnet
\ei

\begin{center}
\includegraphics[max width=\textwidth]{2023_01_13_b07e3b4bcffdaab9792bg-165}
\end{center}


\os{newpage} item[{\rot\bf  Addresses }]~\\[-3ex]

\bi
  \ite Each host and router also needs an IP address at the internet layer to designate its position in the overall Internet
\ei

\begin{center}
\includegraphics[max width=\textwidth]{2023_01_13_b07e3b4bcffdaab9792bg-166}
\end{center}


\os{newpage} item[{\rot\bf  Internet and Data Link Addresses Serve Different Purposes }]~\\[-3ex]

\begin{center}
\includegraphics[max width=\textwidth]{2023_01_13_b07e3b4bcffdaab9792bg-167}
\end{center}

\bi
  \ite IP address
\ei
\bi
 \ite To guide delivery to destination host across the Internet (across multiple networks)
 \ei


\bi
  \ite Subnet Address
\ei
\bi
 \ite To guide delivery between two hosts, two routers, and a host and router within a single subnet
 \ei


\bi
  \ite Same LAN, Frame Relay network, etc.
\ei


\os{newpage} item[{\rot\bf  Analogy }]~\\[-3ex]

\bi
  \ite In company, each person has a company-wide ID number (like IP address)

  \ite In company, person also has a local office number in a building

  \ite Paychecks are made out to ID numbers

  \ite For delivery, also need to know office number

\ei


\os{newpage} item[{\rot\bf  Address Resolution }]~\\[-3ex]

\begin{center}
\includegraphics[max width=\textwidth]{2023_01_13_b07e3b4bcffdaab9792bg-169}
\end{center}

\bi
  \ite Problem
\ei
\bi
 \ite Router knows that destination host is on its subnet based on the IP address of an arriving packet
 \ei

\bi
 \ite Does not know the destination host's subnet address, so cannot deliver the packet across the subnet
 \ei


\begin{center}
\includegraphics[max width=\textwidth]{2023_01_13_b07e3b4bcffdaab9792bg-169(1)}
\end{center}


\os{newpage} item[{\rot\bf  Address Resolution Protocol (ARP) }]~\\[-3ex]

\bi
  \ite Router creates an ARP Request message to be
sent to all hosts on the subnet.

  \ite Address resolution protocol message asks “Who has IP address 128.171.17.13?"

  \ite Passes ARP request to data link layer process for delivery

\ei

\begin{center}
\includegraphics[max width=\textwidth]{2023_01_13_b07e3b4bcffdaab9792bg-170}
\end{center}


\os{newpage} item[{\rot\bf  Address Resolution Protocol (ARP) }]~\\[-3ex]


\os{newpage} item[{\rot\bf  - Data link process of router broadcasts the ARP Request message to all hosts on the subnet. }]~\\[-3ex]


\os{newpage} item[{\rot\bf  o On a LAN, MAC address of 48 ones tells all stations to pay attention to the frame }]~\\[-3ex]

\begin{center}
\includegraphics[max width=\textwidth]{2023_01_13_b07e3b4bcffdaab9792bg-171}
\end{center}


\os{newpage} item[{\rot\bf  Address Resolution Protocol ) }]~\\[-3ex]

\bi
  \ite Host with IP address 128.171.17.13 responds

  \ite Internet process creates an ARP response message

  \ite Contains the destination host's subnet address (48-bit MAC address on a LAN)
\includegraphics[max width=\textwidth, center]{2023_01_13_b07e3b4bcffdaab9792bg-172}

\ei


\item[{\rot\bf  Address Resolution Protocol  }]~\\[-3ex]

\bi
  \ite Router delivers the IP packet to the destination host
\ei
\bi
 \ite Places the IP packet in the subnet frame
 \ei

\bi
 \ite Puts the destination host's subnet address in the destination address field of the frame
 \ei


\begin{center}
\includegraphics[max width=\textwidth]{2023_01_13_b07e3b4bcffdaab9792bg-173}
\end{center}


\os{newpage} item[{\rot\bf  Address Resolution Protocol }]~\\[-3ex]

\bi
  \ite ARP Requests and Responses are sent between thr
internet layer processes on the router and the
destination host
\ei

Router

\begin{center}
\includegraphics[max width=\textwidth]{2023_01_13_b07e3b4bcffdaab9792bg-174}
\end{center}

ARP

Request

\begin{center}
\includegraphics[max width=\textwidth]{2023_01_13_b07e3b4bcffdaab9792bg-174(1)}
\end{center}

ARP

Response Destination Host

Internet
Process

\section{Address Resolution Protocol
 - However, the data link processes deliver these ARP packets
 o Router broadcasts the ARP Request
 o Destination host sends ARP response to the subnet source address found in the broadcast frame}
Router

\begin{center}
\begin{tabular}{|c|}
\hline
Internet \\
Process \\
\hline
Data Link \\
Process \\
\hline
\end{tabular}
\end{center}

Broadcast ARP Request

\begin{center}
\includegraphics[max width=\textwidth]{2023_01_13_b07e3b4bcffdaab9792bg-175}
\end{center}

Direct ARP Response Destination Host

Data Link

Process


\item[{\rot\bf  Direct and Indirect Routing  }]~\\[-3ex]


\os{newpage} item[{\rot\bf  - Direct - when nodes are on the same network }]~\\[-3ex]


\os{newpage} item[{\rot\bf  - Indirect - used when the network numbers of the source and destination do not match }]~\\[-3ex]




\os{newpage} item[{\rot\bf  Transmission Control Protocol (TCP) }]~\\[-3ex]


\os{newpage} item[{\rot\bf  TCP: OvervieW RFCs: 793, 1122, 1323, 2018, 2581 }]~\\[-3ex]

\begin{center}
\includegraphics[max width=\textwidth]{2023_01_13_b07e3b4bcffdaab9792bg-178}
\end{center}

\bi
  \ite point-to-point (unicast):
\ei
\bi
 \ite one sender, one receiver
 \ei


\bi
  \ite connection-oriented:

  \ite handshaking (exchange of control msgs) init's sender, receiver state before data exchange

  \ite State resides only at the END systems - Not a virtual circuit!

  \ite full duplex data:

  \ite bi-directional data flow in same connection (A->B \& B->A in the same connection)

\ei
\bi
 \ite MSS: maximum segment size - reliable, in-order byte steam: o no "message boundaries"
 \ei


\bi
  \ite send \& receive buffers o buffer incoming \& outgoing data

  \ite flow controlled:

\ei
\bi
 \ite sender will not overwhelm receiver
 \ei


\bi
  \ite congestion controlled:
\ei
\bi
 \ite sender will not overwhelm network
 \ei



\os{newpage} item[{\rot\bf  TCP segment structure }]~\\[-3ex]


\os{newpage} item[{\rot\bf  2 bits }]~\\[-3ex]

URG: urgent data (generally not used)

ACK: ACK \# valid

PSH: push data now (generally not used)

RST, SYN, FIN: connection estab (setup, teardown commands)

Internet checksum (as in UDP)
\includegraphics[max width=\textwidth, center]{2023_01_13_b07e3b4bcffdaab9792bg-179}

\begin{center}
\begin{tabular}{|c|c|}
\hline
source port \# & dest port \# \\
\hline
\end{tabular}
\end{center}

by bytes

of data

(not segments!)

$$
\begin{aligned}
& \text { application } \\
& \text { data } \\
& \text { (variable length) }
\end{aligned}
$$

\# bytes

rcvr willing

to accept


\os{newpage} item[{\rot\bf  TCP Connection-oriented demux }]~\\[-3ex]

\bi
  \ite TCP socket identified by 4 -tuple:
\ei
\bi
 \ite source IP address
 \ei

\bi
 \ite source port number
 \ei

\bi
 \ite dest IP address
 \ei

\bi
 \ite dest port number
 \ei



\os{newpage} item[{\rot\bf  - receiving host uses all four values to direct segment to appropriate socket }]~\\[-3ex]


\os{newpage} item[{\rot\bf  TCP Demultiplexing Example }]~\\[-3ex]

\begin{center}
\begin{tabular}{|l|l|l|l|l|l|}
\hline
 &  &  &  \\
\hline
\end{tabular}
\end{center}


\os{newpage} item[{\rot\bf  Typical TCP Transaction }]~\\[-3ex]

\begin{center}
\includegraphics[max width=\textwidth]{2023_01_13_b07e3b4bcffdaab9792bg-182}
\end{center}


\os{newpage} item[{\rot\bf  A TCP Transaction consists of 3 Phases }]~\\[-3ex]

\begin{enumerate}
  \ite Connection Establishment
\end{enumerate}
\bi
 \ite Handshaking between client and server
 \ei



\os{newpage} item[{\rot\bf  Reliable, In-Order Data Exchange }]~\\[-3ex]
\bi
 \ite Recover any lost data through retransmissions and ACKs
 \ei



\os{newpage} item[{\rot\bf  Connection Termination o Closing the connection }]~\\[-3ex]


\os{newpage} item[{\rot\bf  TCP Connection Establishment }]~\\[-3ex]

\bi
  \ite TCP sender, receiver establish "connection" before exchanging data segments

  \ite initialize TCP variables:

\ei
\bi
 \ite seq. \#s
 \ei

\bi
 \ite buffers, flow control info (e.g. RcvWindow)
 \ei


\bi
  \ite client: connection initiator
\ei

Socket clientSocket $=$ new Socket ("hostname", port\#);

\bi
  \ite server: contacted by client Socket connectionSocket $=$ welcomeSocket. accept () ;
\ei


\os{newpage} item[{\rot\bf  Connection Establishment (cont) }]~\\[-3ex]


\os{newpage} item[{\rot\bf  Three way handshake: }]~\\[-3ex]

Step 1: client host sends TCP SYN segment to server
\bi
 \ite specifies a random initial seq \#
 \ei

\bi
 \ite no data
 \ei


Step 2: server host receives SYN, replies with SYNACK segment o server allocates buffers
\bi
 \ite specifies server initial seq. \# Step 3: client receives SYNACK, replies with $\mathrm{ACK}$ segment, which may contain data
 \ei


\begin{center}
\includegraphics[max width=\textwidth]{2023_01_13_b07e3b4bcffdaab9792bg-184}
\end{center}

Host B

Conhection

request $\sqrt{S Y N, S_{e q}=42} \quad$ host ACKs and selects its own initial seq \#


\os{newpage} item[{\rot\bf  Connection Establishment (cont) }]~\\[-3ex]


\os{newpage} item[{\rot\bf  Seq. \#'s: }]~\\[-3ex]
\bi
 \ite byte stream "number" of first byte in segment's data
 \ei



\os{newpage} item[{\rot\bf  ACKs: }]~\\[-3ex]

O seq \# of next byte expected from other side
\bi
 \ite cumulative ACK
 \ei


\begin{center}
\includegraphics[max width=\textwidth]{2023_01_13_b07e3b4bcffdaab9792bg-185}
\end{center}

Host B

Connection

request

\begin{center}
\includegraphics[max width=\textwidth]{2023_01_13_b07e3b4bcffdaab9792bg-185(1)}
\end{center}

host ACKs and selects its own initial seq \#

Three-way handshake


\os{newpage} item[{\rot\bf  TCP Starting Sequence Number Selection }]~\\[-3ex]

\bi
  \ite Why a random starting sequence \#? Why not simply choose o?

  \ite To protect against two incarnations of the same connection reusing the same sequence numbers too soon

  \ite That is, while there is still a chance that a segment from an earlier incarnation of a connection will interfere with a later incarnation of the connection

  \ite How?

  \ite Client machine seq \#o, initiates connection to server with seq \#0.

  \ite Client sends one byte and client machine crashes

  \ite Client reboots and initiates connection again

  \ite Server thinks new incarnation is the same as old connection

\ei


\os{newpage} item[{\rot\bf  TCP Connection Termination }]~\\[-3ex]


\os{newpage} item[{\rot\bf  Closing a connection: }]~\\[-3ex]

client closes socket:

clientSocket. close () ;


\os{newpage} item[{\rot\bf  Step 1: client end system sends TCP FIN control segment to server }]~\\[-3ex]

Step 2: server receives FIN, replies with $\mathrm{ACK}$. Server might send some buffered but not sent data before closing the connection. Server then sends FIN and moves to Closing state.

\begin{center}
\includegraphics[max width=\textwidth]{2023_01_13_b07e3b4bcffdaab9792bg-187}
\end{center}


\os{newpage} item[{\rot\bf  TCP Connection Termination }]~\\[-3ex]

Step 3: client receives FIN, replies with ACK.
\bi
 \ite Enters "timed wait" - will respond with ACK to received FINs
 \ei


Step 4: server, receives ACK. Connection closed.

\bi
  \ite Why wait before closing the connection?

  \ite If the connection were allowed to move to CLOSED state, then another pair of application processes might come along and open the same connection (use the same port \#s) and a delayed FIN from an earlier incarnation would terminate the connection.

\ei

\begin{center}
\includegraphics[max width=\textwidth]{2023_01_13_b07e3b4bcffdaab9792bg-188}
\end{center}


\os{newpage} item[{\rot\bf  TCP State-Transition Diagram }]~\\[-3ex]

\begin{center}
\includegraphics[max width=\textwidth]{2023_01_13_b07e3b4bcffdaab9792bg-189}
\end{center}


\os{newpage} item[{\rot\bf  Typical TCP Client/Server Transitions }]~\\[-3ex]

\begin{center}
\includegraphics[max width=\textwidth]{2023_01_13_b07e3b4bcffdaab9792bg-190}
\end{center}

TCP client

lifecycle TCP server lifecycle

\begin{center}
\includegraphics[max width=\textwidth]{2023_01_13_b07e3b4bcffdaab9792bg-190(1)}
\end{center}


\os{newpage} item[{\rot\bf  How to program using the TCP? }]~\\[-3ex]

\begin{center}
\includegraphics[max width=\textwidth]{2023_01_13_b07e3b4bcffdaab9792bg-191}
\end{center}
\bi
 \ite Programmer's API to the protocol stack
 \ei


$\square$ Typical network app has two pieces: client and server Server: Passive entity. Provides service to clients
\bi
 \ite e.g., Web server responds with the requested Web page
 \ei


$\square$ Client: initiates contact with server ("speaks first")
\bi
 \ite typically requests service from server, e.9., Web Browser
 \ei



\os{newpage} item[{\rot\bf  Socket Creation }]~\\[-3ex]

\bi
  \ite mySock = socket(family, type, protocol);

  \ite UDP/TCP/IP-specific sockets

\ei

\begin{center}
\begin{tabular}{|l|c|c|l|}
\multicolumn{2}{c}{Family} & Type & Protocol \\
\hline
TCP & \multirow{3}{*}{PF\_INET} & SOCK\_STREAM & IPPROTO\_TCP \\
\cline { 3 - 4 }
 &  & SOCK\_DGRAM & IPPROTO\_UDP \\
\hline
UDP &  &  &  \\
\hline
\end{tabular}
\end{center}

\bi
  \ite Socket reference
\ei
\bi
 \ite File (socket) descriptor in UNIX
 \ei

\bi
 \ite Socket handle in WinSock
 \ei



\os{newpage} item[{\rot\bf  TCP Client/Server Interaction }]~\\[-3ex]

Server starts by getting ready to receive client connections...



\begin{enumerate}
  \ite Create a TCP socket

  \ite Establish connection

  \ite Communicate

  \ite Close the connection Server

\end{enumerate}

Create a TCP socket

Assign a port to socket

Set socket to listen

Repeatedly:

a. Accept new connection

b. Communicate

c. Close the connection


\os{newpage} item[{\rot\bf  TCP Client/Server Interaction }]~\\[-3ex]

$/ *$ Create socket for incoming connections */

if $(($ servSock $=$ socket(PF\_INET, SOCK\_STREAM, IPPROTO\_TCP $))<0)$ DieWithError("socket() failed");

Client

\begin{enumerate}
  \ite Create a TCP socket

  \ite Establish connection

  \ite Communicate

  \ite Close the connection Server

\end{enumerate}

Create a TCP socket

Bind socket to a port

Set socket to listen

Repeatedly:

a. Accept new connection

b. Communicate

c. Close the connection


\os{newpage} item[{\rot\bf  TCP Client/Server Interaction }]~\\[-3ex]

echoServAddr.sin\_family $=$ AF\_INET;

echoServAddr.sin\_addr.s\_addr $=$ htonl(INADDR\_ANY);/* Any incoming interface */ echoServAddr.sin\_port $=$ htons(echoServPort);

if (bind(servSock, (struct sockaddr DieWithError("bind() failed");


\os{newpage} item[{\rot\bf  Client }]~\\[-3ex]

\begin{enumerate}
  \ite $\quad$ Create a TCP socket

  \ite Establish connection

  \ite Communicate

  \ite Close the connection /* Internet address family \textit{/ /} Local port */

\end{enumerate}

Server

Create a TCP socket Bind socket to a port

Set socket to listen

Repeatedly:
a. Accept new connection
b. Communicate
c. Close the connection


\item[{\rot\bf  TCP Client/Server Interaction  }]~\\[-3ex]

Client

\begin{enumerate}
  \ite Create a TCP socket

  \ite Establish connection

  \ite Communicate

  \ite Close the connection Server

\end{enumerate}

Create a TCP socket Bind socket to a port

Set socket to listen

Repeatedly:
a. Accept new connection
b. Communicate
c. Close the connection


\item[{\rot\bf  TCP Client/Server Interaction  }]~\\[-3ex]

\textbackslash author\{
for $(i ;) / *$ Run forever $* /$ \textbackslash  \{ \textbackslash  clntLen $=$ sizeof(echoCIntAddr $)$; \textbackslash  if $(($ clntSock=accept(servSock, (struct sockaddr * $)$ \&echoCIntAddr, \&clntLen $))<0)$ \textbackslash  DieWithError("accept() failed");
\}


\os{newpage} item[{\rot\bf  Client }]~\\[-3ex]

\begin{enumerate}
  \ite Create a TCP socket

  \ite Establish connection

  \ite Communicate

  \ite Close the connection

\end{enumerate}


\os{newpage} item[{\rot\bf  Server }]~\\[-3ex]

Create a TCP socket Bind socket to a port Set socket to listen Repeatedly:
a. Accept new connection
b. Communicate
c. Close the connection


\os{newpage} item[{\rot\bf  TCP Client/Server Interaction }]~\\[-3ex]

Server is now blocked waiting for connection from a client



\begin{enumerate}
  \ite Create a TCP socket

  \ite Establish connection

  \ite Communicate

  \ite Close the connection

\end{enumerate}


\os{newpage} item[{\rot\bf  Server }]~\\[-3ex]

Create a TCP socket

Bind socket to a port

Set socket to listen

Repeatedly:
a. Accept new connection
b. Communicate
c. Close the connection


\os{newpage} item[{\rot\bf  TCP Client/Server Interaction }]~\\[-3ex]

Later, a client decides to talk to the server...



\begin{enumerate}
  \ite Create a TCP socket

  \ite Establish connection

  \ite Communicate

  \ite Close the connection Server

\end{enumerate}

Create a TCP socket

Bind socket to a port

Set socket to listen

Repeatedly:
a. Accept new connection
b. Communicate
c. Close the connection


\os{newpage} item[{\rot\bf  TCP Client/Server Interaction }]~\\[-3ex]

/* Create a reliable, stream socket using TCP */

if $(($ sock $=$ socket(PF\_INET, SOCK\_STREAM, IPPROTO\_TCP $))<0)$ DieWithError("socket() failed");


\os{newpage} item[{\rot\bf  Client }]~\\[-3ex]

\begin{enumerate}
  \ite Create a TCP socket

  \ite Establish connection

  \ite Communicate

  \ite Close the connection Server

\end{enumerate}

Create a TCP socket Bind socket to a port

Set socket to listen Repeatedly:

a. Accept new connection

b. Communicate

c. Close the connection


\item[{\rot\bf  TCP Client/Server Interaction  }]~\\[-3ex]

Client

\begin{enumerate}
  \ite Create a TCP socket

  \ite Establish connection

  \ite Communicate

  \ite Close the connection

\end{enumerate}


\os{newpage} item[{\rot\bf  Server }]~\\[-3ex]

Create a TCP socket

Bind socket to a port

Set socket to listen

Repeatedly:
a. Accept new connection
b. Communicate
c. Close the connection


\os{newpage} item[{\rot\bf  TCP Client/Server Interaction }]~\\[-3ex]

echoStringLen $=$ strlen $($ echoString $)$;

/* Determine input length * /

/* Send the string to the server $* /$

if (send(sock, echoString, echoStringLen, 0) $!=$ echoStringLen)

DieWithError("send() sent a different number of bytes than expected");

Client

\begin{enumerate}
  \ite Create a TCP socket

  \ite Establish connection

  \ite Communicate

  \ite Close the connection Server

\end{enumerate}

Create a TCP socket

\begin{enumerate}
  \setcounter{enumi}{1}
  \ite Bind socket to a port

  \ite Set socket to listen

  \ite Repeatedly:

\end{enumerate}

a. Accept new connection

b. Communicate

c. Close the connection


\os{newpage} item[{\rot\bf  TCP Client/Server Interaction }]~\\[-3ex]

/* Receive message from client */

if $(($ recvMsgSize $=\operatorname{recv}($ clntSocket, echoBuffer, RCVBUFSIZE, 0) $)<0)$ DieWithError("recv() failed");

Client

\begin{enumerate}
  \ite Create a TCP socket

  \ite Establish connection

  \ite Communicate

  \ite Close the connection Server

  \ite Create a TCP socket

\end{enumerate}

. Bind socket to a port

\begin{enumerate}
  \setcounter{enumi}{2}
  \ite Set socket to listen

  \ite Repeatedly:
a. Accept new connection
b. Communicate
c. Close the connection

\end{enumerate}


\os{newpage} item[{\rot\bf  TCP Client/Server Interaction }]~\\[-3ex]

close(sock);



\begin{enumerate}
  \ite Create a TCP socket

  \ite Establish connection

  \ite Communicate

  \ite Close the connection

\end{enumerate}


\os{newpage} item[{\rot\bf  close(clntSocket) }]~\\[-3ex]

Server

Create a TCP socket

Bind socket to a port

Set socket to listen

Repeatedly:

a. Accept new connection

b. Communicate

c. Close the connection


\item[{\rot\bf  TCP Tidbits  }]~\\[-3ex]

\bi
  \ite Client knows server address and port

  \ite No correlation between send () and recv ()

\ei

Server

$\operatorname{recv}()->$ "Hello"

$\operatorname{recv}()->$ "Bob"

send("Hi ")

send("Jane")

$\operatorname{recv}()->$ "Hi Jane"


\os{newpage} item[{\rot\bf  Serial Busses }]~\\[-3ex]

\begin{center}
\includegraphics[max width=\textwidth]{2023_01_13_b07e3b4bcffdaab9792bg-206}
\end{center}


\os{newpage} item[{\rot\bf  Serial Buses in Embedded }]~\\[-3ex]

\bi
  \ite UART serial bus for sending debug messages to your
development host

  \ite I2C serial bus for communicating with sensors (e.g., the accelerometer)

\ei


\os{newpage} item[{\rot\bf  - SPI serial bus for communicating with the Bluetooth Low Energy radio }]~\\[-3ex]


\os{newpage} item[{\rot\bf  Serial Interfaces }]~\\[-3ex]

\begin{center}
\includegraphics[max width=\textwidth]{2023_01_13_b07e3b4bcffdaab9792bg-208}
\end{center}


\os{newpage} item[{\rot\bf  Parallel Bus VS Serial Bus }]~\\[-3ex]

\begin{center}
\includegraphics[max width=\textwidth]{2023_01_13_b07e3b4bcffdaab9792bg-209}
\end{center}

\begin{center}
\includegraphics[max width=\textwidth]{2023_01_13_b07e3b4bcffdaab9792bg-209(1)}
\end{center}

\section{Simplistic View of Serial Port Operation Transmitter
 Receiver}
\begin{center}
\includegraphics[max width=\textwidth]{2023_01_13_b07e3b4bcffdaab9792bg-210}
\end{center}

Interrupt raised when Transmitter $(T x)$ is empty $\Rightarrow$ Byte has been transmitted and next byte ready for loading

\begin{center}
\includegraphics[max width=\textwidth]{2023_01_13_b07e3b4bcffdaab9792bg-210(1)}
\end{center}

Interrupt raised when Receiver $(R x)$ is full $\Rightarrow$ Byte has been received and is ready for reading


\os{newpage} item[{\rot\bf  Serial Bus Interface Motivations }]~\\[-3ex]

\bi
  \ite Motivation
\ei
\bi
 \ite Without using a lot of I/O lines
 \ei


$\times \mathrm{I} / \mathrm{O}$ lines require $\mathrm{I} / \mathrm{O}$ pads which cost $\$ \$ \$$ and size

\bi
  \ite I/O lines require PCB area which costs $\$ \$ \$$ and size
\ei
\bi
 \ite Connect different systems together
 \ei


\bi
  \ite Two embedded systems
\ei

$\times$ A desktop and an embedded system
\bi
 \ite Connect different chips together in the same embedded system
 \ei


\bi
  \ite MCU to peripheral
\ei

$\times \mathrm{MCU}$ to $\mathrm{MCU}$
\bi
 \ite Often at relatively low data rates
 \ei


\bi
  \ite But sometimes at higher data rates

  \ite So, what are our options?

\ei
\bi
 \ite Universal Synchronous/Asynchronous Receiver Transmitter
 \ei


\bi
  \ite Also known as USART (pronounced: "you-sart")
\ei


\os{newpage} item[{\rot\bf  Serial Bus Design Space }]~\\[-3ex]

\bi
  \ite Number of wires required?

  \ite Asynchronous or synchronous?

  \ite How fast can it transfer data?

  \ite Can it support more than two endpoints?

  \ite Can it support more than one master (i.e. txn initiator)?

  \ite How do we support flow control?

  \ite How does it handle errors/noise?

  \ite How far can signals travel?

\ei


\os{newpage} item[{\rot\bf  Serial Bus Examples }]~\\[-3ex]

\begin{center}
\includegraphics[max width=\textwidth]{2023_01_13_b07e3b4bcffdaab9792bg-213}
\end{center}


\os{newpage} item[{\rot\bf  UART Uses }]~\\[-3ex]

\bi
  \ite PC serial port is a UART!

  \ite Serializes data to be sent over serial cable - De-serializes received data
\includegraphics[max width=\textwidth, center]{2023_01_13_b07e3b4bcffdaab9792bg-214}

\ei


\os{newpage} item[{\rot\bf  UART Uses }]~\\[-3ex]


\os{newpage} item[{\rot\bf  - Used to be commonly used for internet access }]~\\[-3ex]

\begin{center}
\includegraphics[max width=\textwidth]{2023_01_13_b07e3b4bcffdaab9792bg-215}
\end{center}


\os{newpage} item[{\rot\bf  UART }]~\\[-3ex]

\bi
  \ite Universal Asynchronous Receiver/Transmitter

  \ite Hardware that translates between parallel and serial forms

\ei


\os{newpage} item[{\rot\bf  - Commonly used in conjunction with communication standards such as EIA, RS-232, RS-422 or RS-485 }]~\\[-3ex]

\begin{center}
\includegraphics[max width=\textwidth]{2023_01_13_b07e3b4bcffdaab9792bg-216}
\end{center}


\os{newpage} item[{\rot\bf  Protocol }]~\\[-3ex]

\bi
  \ite Each character is sent as
\ei
\bi
 \ite a logic low start bit
 \ei

\bi
 \ite a configurable number of data bits (usually 7 or 8 , sometimes
 \ei


5.
\bi
 \ite an optional parity bit
 \ei

\bi
 \ite one or more logic high stop bits
 \ei

\bi
 \ite with a particular bit timing ("baud")
 \ei


\begin{center}
\begin{tabular}{|c|c|c|c|c|c|c|c|c|c|}
\hline
Start & Data 0 & Data 1 & Data 2 & Data 3 & Data & Data & Data & Data \\
5 & 6 & 7 & Stop &  &  &  &  &  \\
\hline
\end{tabular}
\end{center}


\os{newpage} item[{\rot\bf  UART Example }]~\\[-3ex]


\os{newpage} item[{\rot\bf  - Send the ASCII letter 'W' (1010111) }]~\\[-3ex]

\begin{center}
\includegraphics[max width=\textwidth]{2023_01_13_b07e3b4bcffdaab9792bg-218}
\end{center}


\os{newpage} item[{\rot\bf  UART Hardware Connection }]~\\[-3ex]

\begin{center}
\includegraphics[max width=\textwidth]{2023_01_13_b07e3b4bcffdaab9792bg-219}
\end{center}


\os{newpage} item[{\rot\bf  UART Character Reception }]~\\[-3ex]

Start bit says a character is coming, receiver resets its timers

\begin{center}
\includegraphics[max width=\textwidth]{2023_01_13_b07e3b4bcffdaab9792bg-220}
\end{center}

Receiver uses a timer (counter) to time when it samples.

Transmission rate (i.e., bit width) must be known!

Slides from BYU CS 224


\os{newpage} item[{\rot\bf  UART Character Reception }]~\\[-3ex]

If receiver samples too quickly, see what happens...

Space

\begin{center}
\includegraphics[max width=\textwidth]{2023_01_13_b07e3b4bcffdaab9792bg-221}
\end{center}


\os{newpage} item[{\rot\bf  UART Character Reception }]~\\[-3ex]

If receiver samples too slowly, see what happens...

Mark

Space

\begin{center}
\includegraphics[max width=\textwidth]{2023_01_13_b07e3b4bcffdaab9792bg-222}
\end{center}

Receiver resynchronizes on every start bit. Only has to be accurate enough to read 9 bits.


\os{newpage} item[{\rot\bf  UART Character Reception }]~\\[-3ex]

\bi
  \ite Receiver also verifies that stop bit is ' 1 '

  \ite If not, reports "framing error" to host system

  \ite New start bit can appear immediately after stop bit - Receiver will resynchronize on each start bit

\ei


\os{newpage} item[{\rot\bf  Let us design a UART transmitter }]~\\[-3ex]

\begin{center}
\includegraphics[max width=\textwidth]{2023_01_13_b07e3b4bcffdaab9792bg-224}
\end{center}


\os{newpage} item[{\rot\bf  Transmitter/System Handshaking }]~\\[-3ex]

\bi
  \ite System asserts Send and holds it high when it wants to send a byte

  \ite UART asserts Busy signal in response

  \ite When UART has finished transfer, UART de-asserts Busy signal

  \ite System de-asserts Send signal
\includegraphics[max width=\textwidth, center]{2023_01_13_b07e3b4bcffdaab9792bg-225}

\ei


\os{newpage} item[{\rot\bf  Transmitter Block Diagram }]~\\[-3ex]

\begin{center}
\includegraphics[max width=\textwidth]{2023_01_13_b07e3b4bcffdaab9792bg-226}
\end{center}

Send

\begin{center}
\begin{tabular}{l}
Busy \\
\hline
ParitySelect \\
\hline
\end{tabular}
\end{center}

ParitySelect

Din $300 \mathrm{HZ}$

Timer

Mod10

Counter State

Machine

Parity

Generator Shift

Register


\os{newpage} item[{\rot\bf  Discussion Questions }]~\\[-3ex]

\bi
  \ite How fast can we run a UART?

  \ite What are the limitations?

  \ite Why do we need start/stop bits?

  \ite How many data bits can be sent?

\ei

19200 baud rate, no parity, 8 data bits, 1 stop bit


\os{newpage} item[{\rot\bf  Serial Peripheral Interconnect (SPI) }]~\\[-3ex]

\bi
  \ite Another kind of serial protocol in embedded systems (proposed by Motorola)

  \ite Four-wire protocol o SCLK - Serial Clock

  \ite MOSI/SIMO - Master Output, Slave Input

  \ite MISO/SOMI - Master Input, Slave Output

  \ite SS - Slave Select

  \ite Single master device and with one or more slave devices

  \ite Higher throughput than I2C and can do "stream transfers"

  \ite No arbitration required

  \ite But

  \ite Requires more pins

  \ite Has no hardware flow control

  \ite No slave acknowledgment (master could be talking to thin air and not even know it)

\ei


\os{newpage} item[{\rot\bf  What is SPI? }]~\\[-3ex]

\bi
  \ite Serial Bus protocol

  \ite Fast, \href{http://Easy.to}{Easy.to} use, Simpleat ..... ° - Everyor'

\ei

\begin{center}
\includegraphics[max width=\textwidth]{2023_01_13_b07e3b4bcffdaab9792bg-229}
\end{center}

\begin{center}
\includegraphics[max width=\textwidth]{2023_01_13_b07e3b4bcffdaab9792bg-229(1)}
\end{center}

Digital Set-Top Box Integrated Controller

POWCNDET

75H6051 03BM 1F11DOORPB KOREA

\begin{center}
\includegraphics[max width=\textwidth]{2023_01_13_b07e3b4bcffdaab9792bg-229(2)}
\end{center}

IBM39 STB02100 PBC 22C




\os{newpage} item[{\rot\bf  SPI Basics }]~\\[-3ex]


\os{newpage} item[{\rot\bf  - A communication protocol using 4 wires Also known as a 4 wire bus }]~\\[-3ex]


\os{newpage} item[{\rot\bf  - Used to communicate across small distances }]~\\[-3ex]


\os{newpage} item[{\rot\bf  - Multiple Slaves, Single Master }]~\\[-3ex]


\os{newpage} item[{\rot\bf  - Synchronized }]~\\[-3ex]


\os{newpage} item[{\rot\bf  SPI Capabilities }]~\\[-3ex]

\bi
  \ite Always Full Duplex

  \ite Communicating in two directions at the same time

  \ite Transmission need not be meaningful

  \ite Multiple Mbps transmission speed

\ei


\os{newpage} item[{\rot\bf  - Transfers data in 4 to 16 bit characters }]~\\[-3ex]

\bi
  \ite Multiple slaves

  \ite Daisy-chaining possible

\ei


\os{newpage} item[{\rot\bf  SPI Protocol }]~\\[-3ex]

\bi
  \ite Wires:
\ei
\bi
 \ite Master Out Slave In (MOSI)
 \ei

\bi
 \ite Master In Slave Out (MISO)
 \ei

\bi
 \ite System Clock (SCLK)
 \ei


\bi
  \ite Slave Select 1...N
\ei

\begin{center}
\includegraphics[max width=\textwidth]{2023_01_13_b07e3b4bcffdaab9792bg-232}
\end{center}

\bi
  \ite Master Generates Clock

  \ite Master Set Slave Select low

  \ite Shift registers shift in and out data

\ei


\os{newpage} item[{\rot\bf  SPI Wires in Detail }]~\\[-3ex]


\os{newpage} item[{\rot\bf  - MOSI - Carries data out of Master to Slave }]~\\[-3ex]

\bi
  \ite MISO - Carries data from Slave to Master - Both signals happen for every transmission

  \ite SS\_BAR - Unique line to select a slave

\ei


\os{newpage} item[{\rot\bf  - SCLK - Master produced clock to synchronize data transfer }]~\\[-3ex]

\begin{center}
\includegraphics[max width=\textwidth]{2023_01_13_b07e3b4bcffdaab9792bg-233}
\end{center}


\os{newpage} item[{\rot\bf  SPI uses a "shift register" model of communications }]~\\[-3ex]


\os{newpage} item[{\rot\bf  Master }]~\\[-3ex]

Slave

\begin{center}
\includegraphics[max width=\textwidth]{2023_01_13_b07e3b4bcffdaab9792bg-234}
\end{center}

Master shifts out data to Slave, and shifts in data from Slave \href{http://upload.wikimedia.org/wikipedia/commons/thumb/b/bb/SPI_8-bit_circular_transfer.svg/400px-SPI_8-bit_circular_transfer.svg.png}{http://upload.wikimedia.org/wikipedia/commons/thumb/b/bb/SPI\_8-bit\_circular\_transfer.svg/400px-SPI\_8-bit\_circular\_transfer.svg.png}


\os{newpage} item[{\rot\bf  SPI Communication }]~\\[-3ex]

\begin{center}
\includegraphics[max width=\textwidth]{2023_01_13_b07e3b4bcffdaab9792bg-235}
\end{center}

SLAVE

SCK

MosI

MISO

Ss

\begin{center}
\includegraphics[max width=\textwidth]{2023_01_13_b07e3b4bcffdaab9792bg-235(1)}
\end{center}

\begin{center}
\begin{tabular}{c|l|l}
SS \\
Slave-Select \\
\end{tabular}
\end{center}


\os{newpage} item[{\rot\bf  SPI clocking: there is no "standard way" }]~\\[-3ex]

\bi
  \ite Four clocking “modes"

  \ite Two phases

  \ite Two polarities

  \ite Master and selected slave must be in the same mode

  \ite During transfers with slaves A and B, Master must

  \ite Configure clock to Slave A's clock mode

  \ite Select Slave A

  \ite Do transfer

  \ite Deselect Slave A

  \ite Configure clock to Slave B's clock mode

  \ite Select Slave B

  \ite Do transfer

  \ite Deselect Slave B

  \ite Master reconfigures clock mode on-the-fly!

\ei


\os{newpage} item[{\rot\bf  SPI timing diagram }]~\\[-3ex]

\begin{center}
\includegraphics[max width=\textwidth]{2023_01_13_b07e3b4bcffdaab9792bg-237}
\end{center}

Timing Diagram - Showing Clock polarities and phases \href{http://www.maxim-ic.com.cn/images/appnotes/3078/3078Fig02.gif}{http://www.maxim-ic.com.cn/images/appnotes/3078/3078Fig02.gif}


\os{newpage} item[{\rot\bf  SPI Pros and Cons }]~\\[-3ex]

\bi
  \ite Pros:

  \ite Fast and easy
Fast for point-to-point connections

Easily allows streaming/Constant data inflow

No addressing/Simple to implement

Everyone supports it

  \ite Cons:

\ei

SS makes multiple slaves very complicated

\bi
  \ite No acknowledgement ability

  \ite No inherent arbitration

\ei

No flow control


\os{newpage} item[{\rot\bf  I2C bus in our projects }]~\\[-3ex]

\bi
  \ite Communication with the accelerometer
\ei

Read from the accelerometer

\bi
  \ite Pros

  \ite Simple wire connection

  \ite Two wires bus that can connect multiple peripherals with the $\mathrm{MCU}$

  \ite Cons

  \ite Complexity is significantly higher

\ei


\os{newpage} item[{\rot\bf  How to operate the camera? }]~\\[-3ex]

Accel

\begin{center}
\begin{tabular}{|l|l|l|}
\hline
MCU \\
\hline
\end{tabular}
\end{center}

\href{https://www.youtube.com/watch?v}{https://www.youtube.com/watch?v} $=e q Z g$

xR6eRjo


\os{newpage} item[{\rot\bf  I2C Details }]~\\[-3ex]

\bi
  \ite Two lines
\ei
\bi
 \ite Serial data line (SDA)
 \ei

\bi
 \ite Serial clock line (SCL)
 \ei



\os{newpage} item[{\rot\bf  - Only two wires for connecting multiple devices }]~\\[-3ex]

\begin{center}
\includegraphics[max width=\textwidth]{2023_01_13_b07e3b4bcffdaab9792bg-241}
\end{center}


\os{newpage} item[{\rot\bf  I2C Details }]~\\[-3ex]


\os{newpage} item[{\rot\bf  - Each I2C device recognized by a unique address }]~\\[-3ex]


\os{newpage} item[{\rot\bf  - Each I2C device can be either a transmitter or receiver }]~\\[-3ex]


\os{newpage} item[{\rot\bf  - I2C devices can be masters or slaves for a data transfer }]~\\[-3ex]
\bi
 \ite Master (usually a microcontroller): Initiates a data transfer on the bus, generates the clock signals to permit that transfer, and terminates the transfer
 \ei

\bi
 \ite Slave: Any device addressed by the master at that time
 \ei



\os{newpage} item[{\rot\bf  Bit Transfer on the $\mathrm{I}^{2} \mathrm{C}$ Bus }]~\\[-3ex]

\bi
  \ite In normal data transfer, the data line only changes state when the clock is low
\ei

$\mathrm{SCL}$ SDA

$\mathrm{SCL}$

\begin{center}
\includegraphics[max width=\textwidth]{2023_01_13_b07e3b4bcffdaab9792bg-243}
\end{center}

Data line stable; Data valid

\begin{center}
\includegraphics[max width=\textwidth]{2023_01_13_b07e3b4bcffdaab9792bg-243(1)}
\end{center}

$\mathrm{d}$


\os{newpage} item[{\rot\bf  Start and Stop Conditions }]~\\[-3ex]

A transition of the data line while the clock line is high is defined as either a start or a stop condition.

Both start and stop conditions are generated by the bus master

The bus is considered busy after a start condition, until a stop condition occurs

\begin{center}
\includegraphics[max width=\textwidth]{2023_01_13_b07e3b4bcffdaab9792bg-244}
\end{center}


\item[{\rot\bf  $\mathrm{I}^{2} \mathrm{C}$ Addressing  }]~\\[-3ex]


\os{newpage} item[{\rot\bf  - Each node has a unique 7 (or 10) bit address }]~\\[-3ex]


\os{newpage} item[{\rot\bf  - Peripherals often have fixed and programmable address portions }]~\\[-3ex]


\os{newpage} item[{\rot\bf  - Addresses starting with o000 or 1111 have special functions:- }]~\\[-3ex]


\bi
  \ite 0000000 Is a General Call Address \textbackslash  o o000001 Is a Null (CBUS) Address \textbackslash  o 1111XXX Address Extension \textbackslash  O 1111111 Address Extension - Next Bytes are the Actual Address
\ei


\os{newpage} item[{\rot\bf  I2C-Connected System }]~\\[-3ex]

\begin{center}
\includegraphics[max width=\textwidth]{2023_01_13_b07e3b4bcffdaab9792bg-246}
\end{center}

Example I2C-connected system with two microcontrollers

(Source: I2C Specification, Philips)


\os{newpage} item[{\rot\bf  Master-Slave Relationships }]~\\[-3ex]

\bi
  \ite Who is the master?
\ei
\bi
 \ite master-transmitters
 \ei

\bi
 \ite master-receivers
 \ei


\bi
  \ite Suppose microcontroller A wants to send information to microcontroller B

  \ite A (master) addresses B (slave)

\ei
\bi
 \ite A (master-transmitter), sends data to B (slave-receiver)
 \ei

\bi
 \ite A terminates the transfer.
 \ei


\bi
  \ite If microcontroller A wants to receive information from microcontroller B

  \ite A (master) addresses microcontroller B (slave)

  \ite A (master-receiver) receives data from B (slave-transmitter)

  \ite A terminates the transfer

  \ite In both cases, the master (microcontroller A) generates the timing and terminates the transfer

\ei


\os{newpage} item[{\rot\bf  Exercise: How fast can I2C run? }]~\\[-3ex]

\begin{center}
\includegraphics[max width=\textwidth]{2023_01_13_b07e3b4bcffdaab9792bg-248}
\end{center}

Capacitor charging in RC circuit

\begin{center}
\includegraphics[max width=\textwidth]{2023_01_13_b07e3b4bcffdaab9792bg-248(1)}
\end{center}

\bi
  \ite How fast can you run it?

  \ite Assumptions

  \ite 0's are driven

  \ite 1's are "pulled up"

  \ite Some working figures

  \ite $\mathrm{R}_{\mathrm{p}}=10 \mathrm{k} \Omega$

  \ite $\mathrm{C}_{\text {cap }}=100 \mathrm{pF}$

  \ite $V_{D D}=5 V$

  \ite $ V_{in_high }=3.5 ~V $

  \ite Recall for RC circuit

\ei

$-\mathrm{V}_{\text {cap }}(\mathrm{t})=\mathrm{V}_{\mathrm{DD}}\left(1-\mathrm{e}^{-\mathrm{t} / \mathrm{\tau}}\right)$

\bi
  \ite Where $\tau=R C$
\ei


\os{newpage} item[{\rot\bf  Exercise: Bus bit rate vs Useful data rate }]~\\[-3ex]

\bi
  \ite An I2C "transactions" involves the following bits

  \ite $\langle\mathrm{S}\rangle\langle\mathrm{A} 6: \mathrm{A} 0\rangle\langle\mathrm{R} / \mathrm{W}\rangle\langle\mathrm{A}\rangle\langle\mathrm{D} 7: \mathrm{D} 0\rangle\langle\mathrm{A}\rangle\langle\mathrm{F}\rangle$

  \ite Which of these actually carries useful data?

  \ite $\langle\mathrm{S}\rangle\langle\mathrm{A} 6: \mathrm{A} 0\rangle\langle\mathrm{R} / \mathrm{W}\rangle\langle\mathrm{A}\rangle\langle\mathrm{D} 7: \mathrm{D} 0\rangle\langle\mathrm{A}\rangle\langle\mathrm{F}\rangle$

  \ite So, if a bus runs at $400 \mathrm{kHz}$

  \ite What is the clock period?

  \ite What is the data throughput (i.e. data-bits/second)?

  \ite What is the bus "efficiency"?

\ei

\end{description}
